\documentclass[12pt,openany]{book}
\usepackage{amsmath}
\usepackage{amssymb}
\usepackage{amsthm}
\usepackage{mathtools}
\usepackage[hidelinks]{hyperref}
\usepackage{booktabs}
\usepackage{geometry}
\usepackage{fancyhdr}
\usepackage{setspace}
\usepackage{microtype}
\usepackage{xcolor}
\usepackage{enumitem}
\usepackage{titlesec}
\usepackage{tocloft}
\usepackage{array}
\usepackage{tabularx}
\usepackage{float}
\usepackage[all]{hypcap}    % fix hyperref anchor duplication for [H] floats
\usepackage{thmtools}

% ============================================================
% CUSTOM MACROS
% ============================================================
\newcommand{\R}{\mathbb{R}}

% ============================================================
% GEOMETRY AND LAYOUT SETTINGS
% ============================================================
\geometry{margin=1in, includefoot}
\microtypesetup{expansion=true, final}
\singlespacing
\setlength{\parindent}{0.5in}
\setlength{\parskip}{0.5ex plus 0.2ex minus 0.1ex}
\setlength{\emergencystretch}{1.5em}
\tolerance=9999
\hbadness=10000
\vbadness=10000
\hyphenpenalty=10000
\exhyphenpenalty=10000
\raggedbottom
\sloppy

% ============================================================
% SECTION FORMATTING
% ============================================================
\titleformat{\chapter}[display]
    {\normalfont\Huge\bfseries}
    {\chaptertitlename\ \thechapter}{20pt}{\Huge}
\titlespacing*{\chapter}{0pt}{50pt}{40pt}
\titleformat{\section}{\normalfont\Large\bfseries}{\thesection}{1em}{}
\titlespacing*{\section}{0pt}{3.5ex plus 1ex minus 0.2ex}{2.3ex plus 0.2ex}
\titleformat{\subsection}{\normalfont\large\bfseries}{\thesubsection}{1em}{}
\titlespacing*{\subsection}{0pt}{3.25ex plus 1ex minus 0.2ex}{1.5ex plus 0.2ex}
\titleformat{\subsubsection}{\normalfont\normalsize\bfseries}{\thesubsubsection}{1em}{}
\titlespacing*{\subsubsection}{0pt}{3ex plus 0.5ex minus 0.2ex}{1.5ex plus 0.2ex}

% ============================================================
% HEADER AND FOOTER
% ============================================================
\pagestyle{fancy}
\fancyhf{}
\fancyhead[LE]{\thepage}
\fancyhead[RO]{\thepage}
\fancyhead[RE]{\nouppercase{\leftmark}}
\fancyhead[LO]{\nouppercase{\rightmark}}
\renewcommand{\headrulewidth}{0.4pt}
\renewcommand{\footrulewidth}{0pt}
\setlength{\headheight}{14.5pt}

% ============================================================
% EQUATION SETTINGS
% ============================================================
\allowdisplaybreaks[3]
\setlength{\abovedisplayskip}{6pt plus 2pt minus 3pt}
\setlength{\belowdisplayskip}{6pt plus 2pt minus 3pt}
\setlength{\abovedisplayshortskip}{3pt plus 1pt minus 1pt}
\setlength{\belowdisplayshortskip}{3pt plus 1pt minus 1pt}
\numberwithin{equation}{chapter}

% ============================================================
% THEOREM ENVIRONMENTS (shared counter, per math_proof_standards.txt Section 2.2.4)
% ============================================================
\theoremstyle{plain}
\newtheorem{theorem}{Theorem}[chapter]
\newtheorem{lemma}[theorem]{Lemma}
\newtheorem{corollary}[theorem]{Corollary}
\newtheorem{proposition}[theorem]{Proposition}
\newtheorem{conjecture}[theorem]{Conjecture}
\theoremstyle{definition}
\newtheorem{definition}[theorem]{Definition}
\newtheorem{axiom}[theorem]{Axiom}
\newtheorem{assumption}[theorem]{Assumption}
\theoremstyle{remark}
\newtheorem{remark}[theorem]{Remark}

\newenvironment{abstract}{%
  \chapter*{Abstract}
  \begin{quote}\small
}{%
  \end{quote}
}
\newenvironment{keywords}{%
  \paragraph*{Keywords:}
}{%
}
\newenvironment{executivesummary}{%
  \chapter*{Executive Summary}
}{%
}
\newenvironment{keyfindings}{%
  \chapter*{Key Findings}
}{%
}
% ============================================================
% SECTION FORMATTING
% ============================================================
\titleformat{\chapter}[display]
    {\normalfont\Huge\bfseries}
    {\chaptertitlename\ \thechapter}{20pt}{\Huge}
\titlespacing*{\chapter}{0pt}{50pt}{40pt}
\titleformat{\section}{\normalfont\Large\bfseries}{\thesection}{1em}{}
\titlespacing*{\section}{0pt}{3.5ex plus 1ex minus 0.2ex}{2.3ex plus 0.2ex}
\titleformat{\subsection}{\normalfont\large\bfseries}{\thesubsection}{1em}{}
\titlespacing*{\subsection}{0pt}{3.25ex plus 1ex minus 0.2ex}{1.5ex plus 0.2ex}
\titleformat{\subsubsection}{\normalfont\normalsize\bfseries}{\thesubsubsection}{1em}{}
\titlespacing*{\subsubsection}{0pt}{3ex plus 0.5ex minus 0.2ex}{1.5ex plus 0.2ex}

% ============================================================
% TABLE OF CONTENTS FORMATTING
% ============================================================
\renewcommand{\cftchapleader}{\cftdotfill{\cftdotsep}}
\setlength{\cftchapindent}{0em}
\setlength{\cftsecindent}{1.5em}

% ============================================================
% HEADER AND FOOTER
% ============================================================
\pagestyle{fancy}
\fancyhf{}
\fancyhead[LE]{\thepage}
\fancyhead[RO]{\thepage}
\fancyhead[RE]{\nouppercase{\leftmark}}
\fancyhead[LO]{\nouppercase{\rightmark}}
\renewcommand{\headrulewidth}{0.4pt}
\renewcommand{\footrulewidth}{0pt}
\setlength{\headheight}{14.5pt}

% ============================================================
% EQUATION SETTINGS
% ============================================================
\allowdisplaybreaks[3]
\setlength{\abovedisplayskip}{6pt plus 2pt minus 3pt}
\setlength{\belowdisplayskip}{6pt plus 2pt minus 3pt}
\setlength{\abovedisplayshortskip}{3pt plus 1pt minus 1pt}
\setlength{\belowdisplayshortskip}{3pt plus 1pt minus 1pt}
\numberwithin{equation}{chapter}

% ============================================================
% THEOREM ENVIRONMENTS
% ============================================================
\theoremstyle{plain}


% Unique hyperref anchors for every theorem-like environment sharing the
% chapter-numbered theorem counter (prevents duplicate-destination warnings
% for the first theorem-like environment of each chapter)
\providecommand{\theHtheorem}{}
\renewcommand{\theHtheorem}{\thechapter.\arabic{theorem}}
\providecommand{\theHlemma}{}
\renewcommand{\theHlemma}{\thechapter.\arabic{theorem}}
\providecommand{\theHcorollary}{}
\renewcommand{\theHcorollary}{\thechapter.\arabic{theorem}}
\providecommand{\theHproposition}{}
\renewcommand{\theHproposition}{\thechapter.\arabic{theorem}}
\providecommand{\theHconjecture}{}
\renewcommand{\theHconjecture}{\thechapter.\arabic{theorem}}
\providecommand{\theHdefinition}{}
\renewcommand{\theHdefinition}{\thechapter.\arabic{theorem}}
\providecommand{\theHaxiom}{}
\renewcommand{\theHaxiom}{\thechapter.\arabic{theorem}}
\providecommand{\theHassumption}{}
\renewcommand{\theHassumption}{\thechapter.\arabic{theorem}}
\providecommand{\theHremark}{}
\renewcommand{\theHremark}{\thechapter.\arabic{theorem}}
\makeatletter
\providecommand*{\toclevel@conjecture}{0}
\makeatother


% Real and Ideal environments
% Note: \Real and \Ideal are math commands defined in CUSTOM COMMANDS below.
% The block environments are named RealRegime and IdealRegime to avoid conflicts.

% NOTE: All framework macros (\V, \Ideal, \Real, \Naive, \Oracle, \Frontier,
% \Gap, \Reward, \Success, \Perf) are defined once in the CUSTOM COMMANDS
% block below.  Any new symbol must be added there before first use.





% ============================================================
% ENUMERATE-ONLY DISCIPLINE (itemize forbidden)
% ============================================================
\setlist[enumerate]{leftmargin=*, partopsep=0.1ex}

% ============================================================
% PREVENT WIDOWS AND ORPHANS
% ============================================================
\clubpenalty=10000
\widowpenalty=10000
\displaywidowpenalty=10000
\renewcommand{\arraystretch}{1.2}
\setlength{\tabcolsep}{4pt}

% hyperref is loaded with the hidelinks option in the package block above;
% no additional \hypersetup call is required here.

% ============================================================
% MATH OPERATORS
% ============================================================
\DeclareMathOperator{\dom}{dom}
\DeclareMathOperator{\ran}{ran}
\DeclareMathOperator{\CVaR}{CVaR}
\DeclareMathOperator{\VaR}{VaR}
\DeclareMathOperator{\Var}{Var}
\DeclareMathOperator{\Cov}{Cov}

% ============================================================
% CUSTOM COMMANDS
% ============================================================
\newcommand{\Ideal}{\mathcal{E}_{I}}
\newcommand{\Real}{\mathcal{E}_{R}}
\newcommand{\Naive}{N}
\newcommand{\Oracle}{O}
\newcommand{\Frontier}{F}
\newcommand{\Gap}{\mathcal{G}}
\newcommand{\V}{V}
\newcommand{\Reward}{\mathrm{Rew}}
\newcommand{\Success}{\mathrm{Succ}}
\newcommand{\Perf}{\mathrm{Perf}}
\newcommand{\E}{\mathbb{E}}
\newcommand{\N}{\mathbb{N}}
\newcommand{\calR}{\mathcal{R}}
\newcommand{\calO}{\mathcal{O}}
\newcommand{\calS}{\mathcal{S}}
\newcommand{\calI}{\mathcal{I}}
\newcommand{\calU}{\mathcal{U}}
\newcommand{\calL}{\mathcal{L}}
\newcommand{\calT}{\mathcal{T}}
\newcommand{\bfw}{\mathbf{w}}
\newcommand{\bfx}{\mathbf{x}}
\newcommand{\bfmu}{\boldsymbol{\mu}}
\newcommand{\bfSigma}{\boldsymbol{\Sigma}}
\newcommand{\bfA}{\mathbf{A}}
\newcommand{\bfQ}{\mathbf{Q}}
\newcommand{\MaxDD}{\mathrm{MaxDD}}
\newcommand{\QPC}{\mathrm{QPC}}



\title{Ideal--Real Regime Separation, Na\"{i}ve--Oracle Duality, and\\
Frontier-Model Evaluation: A Formal Proof Framework\\begin{equation}0.6em]
\large Mathematical Foundations of the Multi-Agent\\
Quantitative-Finance Platform}
\author{Hadrian Hu}
\date{}

\begin{document}

\maketitle
\tableofcontents
\listoftheorems[ignoreall,show={theorem,lemma,proposition,corollary,definition,axiom,assumption,conjecture}]

% ==============================================================
\begin{abstract}
This monograph establishes the formal mathematical foundations underpinning
a multi-agent, cross-market quantitative-finance platform.  It covers
uncertainty-aware signal aggregation, Kelly-criterion position sizing,
risk-adjusted performance metrics, Bayesian market-regime inference,
portfolio optimisation under uncertainty, Conditional Value-at-Risk,
dynamic hedge-ratio derivation, the temporal-leakage invariant, operator
composition algebra, four-domain asset partition, dependency-graph
acyclicity, idempotency and referential transparency, the risk-adjusted
objective function, type-safe language-layer assignment, Hidden~Markov
regime detection, and epistemic--aleatoric uncertainty decomposition.
All results are stated as formal definitions, axioms, and theorems with
complete proofs or proof sketches.
\end{abstract}

\begin{keywords}
quantitative finance; multi-agent systems; operator composition; market-regime
classification; Bayesian inference; portfolio optimisation; Kelly criterion;
Conditional Value-at-Risk; hedge ratio; temporal leakage; uncertainty
quantification; Hidden Markov Model; epistemic uncertainty; aleatoric uncertainty;
Sharpe ratio; Sortino ratio; Calmar ratio; SOLID principles; polyglot architecture
\end{keywords}

% ==============================================================
\chapter{Uncertainty-Aware Signal Aggregation}
\label{ch:signal-agg}
% ==============================================================

\section{Setting}

We consider a system of $N$ independent signal operators each producing a
scalar signal and an associated confidence.

\begin{definition}[Signal and Confidence]
\label{def:signal-confidence}
A \emph{signal--confidence pair} is a tuple $(s_i, c_i)$ where
$s_i \in [-1, 1]$ is the directional signal produced by operator~$i$ and
$c_i \in (0, 1]$ is the operator's self-reported confidence weight,
for $i = 1, \ldots, N$.
\end{definition}

\section{Aggregate Signal}

\begin{definition}[Confidence-Weighted Aggregate Signal]
\label{def:agg-signal}
Given signal--confidence pairs $\{(s_i, c_i)\}_{i=1}^{N}$, the
\emph{confidence-weighted aggregate signal} is
\begin{equation}
  S \;=\; \frac{\displaystyle\sum_{i=1}^{N} c_i\,s_i}
               {\displaystyle\sum_{i=1}^{N} c_i}.
  \label{eq:agg-signal}
\end{equation}
The \emph{aggregate uncertainty} is
\begin{equation}
  U \;=\; 1 - \frac{\displaystyle N\sum_{i=1}^{N} c_i^{2}}
                   {\Bigl(\displaystyle\sum_{i=1}^{N} c_i\Bigr)^{2}}.
  \label{eq:agg-uncertainty}
\end{equation}
\end{definition}

\begin{theorem}[Boundedness of the Aggregate Signal]
\label{thm:signal-bounded}
The aggregate signal $S$ defined in~\eqref{eq:agg-signal} satisfies
$S \in [-1, 1]$.
\end{theorem}

\begin{proof}
We establish the result by validating all assumptions, enumerating every
case, and justifying each algebraic step without exception.

\begin{enumerate}
  \item \textbf{Assumption validation --- positivity of weights.}
        By hypothesis, $c_i > 0$ for all $i \in \{1,\dots,N\}$.
        Consequently,
        \begin{equation}
          C \;\coloneqq\; \sum_{i=1}^{N} c_i \;>\; 0,
          \label{eq:C-pos}
        \end{equation}
        so the denominator of $S$ in~\eqref{eq:agg-signal} is strictly
        positive and the quotient is well-defined.

  \item \textbf{Assumption validation --- boundedness of individual signals.}
        By hypothesis, $s_i \in [-1, 1]$ for all $i$, which is equivalent
        to $|s_i| \leq 1$.

  \item \textbf{Assumption validation --- non-emptiness of the index set.}
        We require $N \geq 1$; the case $N = 0$ renders the
        sum~\eqref{eq:agg-signal} empty and the formula undefined.
        The theorem is stated for $N \geq 1$ signals, so this edge case
        is excluded by hypothesis.

  \item \textbf{Bounding each weighted term.}
        For every $i \in \{1,\dots,N\}$, since $c_i > 0$ and
        $|s_i| \leq 1$,
        \begin{align}
          c_i\,s_i &\;\leq\; c_i\,|s_i| \;\leq\; c_i \cdot 1 \;=\; c_i,
          \label{eq:upper-term}\\
          c_i\,s_i &\;\geq\; -c_i\,|s_i| \;\geq\; -c_i.
          \label{eq:lower-term}
        \end{align}

  \item \textbf{Summing the upper bound.}
        Summing~\eqref{eq:upper-term} over $i = 1,\dots,N$:
        \begin{equation}
          \sum_{i=1}^{N} c_i\,s_i
            \;\leq\; \sum_{i=1}^{N} c_i \;=\; C.
          \label{eq:sum-upper}
        \end{equation}

  \item \textbf{Summing the lower bound.}
        Summing~\eqref{eq:lower-term} over $i = 1,\dots,N$:
        \begin{equation}
          \sum_{i=1}^{N} c_i\,s_i
            \;\geq\; -\sum_{i=1}^{N} c_i \;=\; -C.
          \label{eq:sum-lower}
        \end{equation}

  \item \textbf{Combining into an absolute-value bound.}
        From~\eqref{eq:sum-upper} and~\eqref{eq:sum-lower},
        \begin{equation}
          -C \;\leq\; \sum_{i=1}^{N} c_i\,s_i \;\leq\; C.
          \label{eq:combined}
        \end{equation}
        By the standard equivalence $|x| \leq M \Leftrightarrow -M \leq x
        \leq M$ (for $M \geq 0$),~\eqref{eq:combined} is equivalent to
        \begin{equation}
          \left|\sum_{i=1}^{N} c_i\,s_i\right| \;\leq\; C.
          \label{eq:abs-num}
        \end{equation}

  \item \textbf{Dividing by the positive denominator.}
        Since $C > 0$ by~\eqref{eq:C-pos}, dividing both sides
        of~\eqref{eq:abs-num} by $C$ preserves the inequality direction:
        \begin{align}
          \frac{\left|\displaystyle\sum_{i=1}^{N} c_i\,s_i\right|}{C}
            &\;\leq\; \frac{C}{C} \;=\; 1.
          \label{eq:div-C}
        \end{align}

  \item \textbf{Applying the absolute-value quotient rule.}
        For any real numerator $A$ and positive denominator $C > 0$,
        \begin{equation}
          \left|\frac{A}{C}\right| \;=\; \frac{|A|}{C}.
          \label{eq:abs-quotient}
        \end{equation}
        Setting $A = \sum_{i=1}^{N} c_i\,s_i$ and applying
        \eqref{eq:abs-quotient} to the definition~\eqref{eq:agg-signal}:
        \begin{align}
          |S|
            &\;=\; \left|\frac{\displaystyle\sum_{i=1}^{N} c_i\,s_i}{C}
                   \right|
             \;=\; \frac{\left|\displaystyle\sum_{i=1}^{N} c_i\,s_i\right|}
                        {C}
             \;\leq\; 1.
          \label{eq:abs-S}
        \end{align}

  \item \textbf{Boundary case --- all signals at the upper extremity.}
        If $s_i = 1$ for all $i$, then
        $S = C / C = 1$, so the upper bound is achieved.

  \item \textbf{Boundary case --- all signals at the lower extremity.}
        If $s_i = -1$ for all $i$, then
        $S = -C / C = -1$, so the lower bound is achieved.

  \item \textbf{Boundary case --- mixed signals.}
        For any mixture of $s_i \in [-1, 1]$, Steps~4--9 apply without
        modification; no special case arises because the triangle inequality
        and the positivity of $C$ are preserved for all finite, non-empty
        index sets.

  \item \textbf{Conclusion.}
        From~\eqref{eq:abs-S}, $|S| \leq 1$, which by the equivalence
        $|x| \leq 1 \Leftrightarrow x \in [-1, 1]$ gives
        \begin{equation}
          S \;\in\; [-1,\, 1]. \qed
        \end{equation}
\end{enumerate}
\end{proof}

\begin{remark}
A \texttt{ConfidenceAggregationOperator} must preserve the bound
$S \in [-1, 1]$ and must not silently drop low-confidence signals without
recording the provenance of the omission.
\end{remark}

% ==============================================================
\chapter{Kelly Criterion and Position Sizing}
\label{ch:kelly}
% ==============================================================

\section{Problem Statement}

Given win probability $p$, loss probability $q = 1 - p$, and a
win/loss payoff ratio $b > 0$ (dollars won per dollar bet on a win), we
seek the fraction $f^{*}$ of capital to wager that maximises long-run
logarithmic wealth growth.

\begin{definition}[Logarithmic Growth Rate]
\label{def:log-growth}
For a bet fraction $f \in (0, 1)$, the \emph{expected logarithmic growth
rate} per trial is
\begin{equation}
  G(f) \;=\; p\ln(1 + bf) + q\ln(1 - f).
  \label{eq:log-growth}
\end{equation}
\end{definition}

\section{Derivation of the Kelly Fraction}

% --------------------------------------------------------------
\subsection*{Standing Assumptions for Chapter~\ref{ch:kelly}}
\label{sec:kelly-assumptions}
% --------------------------------------------------------------

The following assumptions are assumed to hold throughout this chapter
and are validated prior to every theorem.

\begin{enumerate}[label=\textbf{A\arabic*.}, leftmargin=*, widest=A9]
  \item \label{asm:prob-valid}
        $p \in (0, 1)$ and $q \coloneqq 1 - p \in (0, 1)$, so both
        outcomes occur with strictly positive probability.  In
        particular, the degenerate cases $p = 0$ and $p = 1$ are
        excluded.
  \item \label{asm:payoff-positive}
        $b > 0$, i.e.\ the payoff ratio is strictly positive.  A zero
        or negative payoff ratio would render the betting problem
        ill-posed.
  \item \label{asm:fraction-domain}
        The bet fraction satisfies $f \in (0, 1)$, so neither full
        abstention ($f = 0$) nor ruinous over-betting ($f \geq 1$) is
        permitted.  On the boundary $f = 1$ the term $\ln(1 - f)$
        diverges to $-\infty$, so $f = 1$ is excluded from the
        feasible domain.
  \item \label{asm:kelly-positive}
        We assume throughout that $pb > q$, i.e.\ $f^{*} > 0$, which
        is the necessary and sufficient condition for the bet to have
        positive expected log-growth.  The complementary case
        $pb \leq q$ is addressed in Remark~\ref{rem:negative-kelly}.
\end{enumerate}

% --------------------------------------------------------------
\subsection*{Domain Analysis of $G$}
\label{sec:kelly-domain}
% --------------------------------------------------------------

\begin{proposition}[Regularity of $G$]
\label{prop:G-regularity}
Under Assumptions \ref{asm:prob-valid}--\ref{asm:fraction-domain},
the function $G : (0,1) \to \mathbb{R}$ defined by
\eqref{eq:log-growth} is:
\begin{enumerate}[label=(\roman*)]
  \item well-defined and finite on all of $(0, 1)$;
  \item infinitely differentiable ($C^{\infty}$) on $(0, 1)$;
  \item strictly concave on $(0, 1)$.
\end{enumerate}
\end{proposition}

\begin{proof}
We verify each of the three claimed properties in turn, validating
every relevant assumption at each step.

\begin{enumerate}[label=\textbf{Step \arabic*.}, leftmargin=*, widest=Step 9]

  % -----------------------------------------------------------
  \item \textbf{Validation of Assumption~\ref{asm:prob-valid}
        (probability validity).}
        By Assumption~\ref{asm:prob-valid}, $p \in (0,1)$ and
        $q = 1 - p$, so:
        \begin{align}
          p   &> 0,                        \label{eq:p-pos} \\
          q   &= 1 - p \;\in\; (0,1),     \label{eq:q-pos} \\
          p+q &= 1.                        \label{eq:pq-sum}
        \end{align}
        In particular $p > 0$ and $q > 0$, which are used
        repeatedly below.

  % -----------------------------------------------------------
  \item \textbf{Validation of Assumption~\ref{asm:odds-valid}
        (odds positivity).}
        By Assumption~\ref{asm:odds-valid}, $b > 0$.
        Combined with \eqref{eq:p-pos} this gives:
        \begin{align}
          pb &> 0.  \label{eq:pb-pos}
        \end{align}

  % -----------------------------------------------------------
  \item \textbf{Validation of Assumption~\ref{asm:fraction-domain}
        (feasible domain).}
        The fraction $f$ is restricted to the open interval
        $(0,1)$.  We note:
        \begin{align}
          f   &\in (0,1)
              \;\Longrightarrow\;
              f > 0 \text{ and } f < 1.
              \label{eq:f-domain}
        \end{align}
        Both strict inequalities hold simultaneously throughout
        the domain; no boundary point is included.

  % -----------------------------------------------------------
  \item \textbf{Well-definedness of $G$ on $(0,1)$.}
        Recall the definition
        \begin{align}
          G(f) &= p\ln(1 + bf) + q\ln(1 - f).
                 \label{eq:G-def-repeated}
        \end{align}
        The natural logarithm $\ln(\,\cdot\,)$ is defined and
        finite on $(0,\infty)$.  We must verify that both
        arguments lie in $(0,\infty)$ for every $f\in(0,1)$.

        \begin{enumerate}[label=(\alph*)]
          \item \emph{First argument.}
                Since $b > 0$ (Step~2) and $f > 0$ (Step~3):
                \begin{align}
                  bf     &> 0,                     \label{eq:bf-pos}\\
                  1 + bf &> 1 > 0.                 \label{eq:arg1-pos}
                \end{align}
                Moreover, $f < 1$ gives an explicit upper bound:
                \begin{align}
                  1 + bf &< 1 + b.                 \label{eq:arg1-upper}
                \end{align}
                Hence $1 + bf \in (1,\, 1+b) \subset (0,\infty)$.

          \item \emph{Second argument.}
                Since $f < 1$ (Step~3):
                \begin{align}
                  1 - f &\in (0,1) \subset (0,\infty). \label{eq:arg2-pos}
                \end{align}
        \end{enumerate}

        Both arguments are strictly positive, so both logarithms
        are finite.  Since $p, q > 0$ (Steps~1--2), the linear
        combination \eqref{eq:G-def-repeated} is finite and
        well-defined for every $f \in (0,1)$. \hfill$\checkmark$

  % -----------------------------------------------------------
  \item \textbf{Smoothness ($G \in C^{\infty}(0,1)$).}
        We identify $G$ as a composition and linear combination
        of elementary functions:
        \begin{enumerate}[label=(\alph*)]
          \item The affine maps $f \mapsto 1+bf$ and
                $f \mapsto 1-f$ are polynomials, hence
                $C^{\infty}(\mathbb{R})$.
          \item From Steps~4(a)--4(b), both affine maps take
                values in $(1, 1+b)$ and $(0,1)$ respectively
                on $(0,1)$; in particular their images are
                contained in $(0,\infty)$.
          \item The function $\ln : (0,\infty)\to\mathbb{R}$
                is $C^{\infty}$ on its domain.
          \item A composition of $C^{\infty}$ functions whose
                intermediate values remain in the domain of the
                outer function is itself $C^{\infty}$.
        \end{enumerate}
        Therefore each of $f \mapsto \ln(1+bf)$ and
        $f \mapsto \ln(1-f)$ is $C^{\infty}$ on $(0,1)$, and
        a finite linear combination with constant coefficients
        preserves this regularity.  Hence:
        \begin{align}
          G \in C^{\infty}(0,1). \label{eq:G-smooth}
        \end{align}
        \hfill$\checkmark$

  % -----------------------------------------------------------
  \item \textbf{First derivative of $G$.}
        Differentiating \eqref{eq:G-def-repeated} termwise
        (justified by $C^{\infty}$ regularity from Step~5):
        \begin{align}
          \frac{d}{df}\bigl[p\ln(1+bf)\bigr]
            &= \frac{pb}{1+bf},
               \label{eq:dG-term1}\\[4pt]
          \frac{d}{df}\bigl[q\ln(1-f)\bigr]
            &= \frac{q\cdot(-1)}{1-f}
             = -\frac{q}{1-f}.
               \label{eq:dG-term2}
        \end{align}
        Summing \eqref{eq:dG-term1} and \eqref{eq:dG-term2}:
        \begin{align}
          G'(f) &= \frac{pb}{1+bf} - \frac{q}{1-f}.
                  \label{eq:G-prime}
        \end{align}
        Both denominators are strictly positive on $(0,1)$ by
        Steps~4(a)--4(b), so $G'(f)$ is well-defined and
        finite on $(0,1)$.

  % -----------------------------------------------------------
  \item \textbf{Second derivative of $G$.}
        Differentiating \eqref{eq:G-prime} termwise:
        \begin{align}
          \frac{d}{df}\!\left[\frac{pb}{1+bf}\right]
            &= pb \cdot \frac{-b}{(1+bf)^{2}}
             = -\frac{pb^{2}}{(1+bf)^{2}},
               \label{eq:d2G-term1}\\[4pt]
          \frac{d}{df}\!\left[-\frac{q}{1-f}\right]
            &= -q \cdot \frac{1}{(1-f)^{2}}
             = -\frac{q}{(1-f)^{2}}.
               \label{eq:d2G-term2}
        \end{align}
        Summing \eqref{eq:d2G-term1} and \eqref{eq:d2G-term2}:
        \begin{align}
          G''(f) &= -\frac{pb^{2}}{(1+bf)^{2}}
                    -\frac{q}{(1-f)^{2}}.
                    \label{eq:G-double-prime}
        \end{align}

  % -----------------------------------------------------------
  \item \textbf{Sign analysis of $G''(f)$ and strict concavity.}
        We analyse each term in \eqref{eq:G-double-prime}
        separately for all $f \in (0,1)$.

        \begin{enumerate}[label=(\alph*)]
          \item \emph{First term.}
                \begin{align}
                  p          &> 0
                    && \text{(Step~1, \eqref{eq:p-pos})},
                    \label{eq:cc-p}\\
                  b^{2}      &> 0
                    && \text{(Step~2, } b>0\text{)},
                    \label{eq:cc-b2}\\
                  (1+bf)^{2} &> 0
                    && \text{(Step~4(a), \eqref{eq:arg1-pos})},
                    \label{eq:cc-denom1}\\
                  \Rightarrow\quad
                  -\frac{pb^{2}}{(1+bf)^{2}} &< 0.
                    \label{eq:cc-t1-neg}
                \end{align}

          \item \emph{Second term.}
                \begin{align}
                  q          &> 0
                    && \text{(Step~1, \eqref{eq:q-pos})},
                    \label{eq:cc-q}\\
                  (1-f)^{2}  &> 0
                    && \text{(Step~4(b), \eqref{eq:arg2-pos})},
                    \label{eq:cc-denom2}\\
                  \Rightarrow\quad
                  -\frac{q}{(1-f)^{2}} &< 0.
                    \label{eq:cc-t2-neg}
                \end{align}
        \end{enumerate}

        Since both terms are strictly negative, their sum
        satisfies:
        \begin{align}
          G''(f)
            &= \underbrace{-\frac{pb^{2}}{(1+bf)^{2}}}_{<\,0}
             + \underbrace{\left(-\frac{q}{(1-f)^{2}}\right)}_{<\,0}
             < 0
             \qquad \forall\, f \in (0,1).
               \label{eq:G-pp-neg}
        \end{align}
        A function that is twice differentiable on an open
        interval with strictly negative second derivative
        throughout that interval is \emph{strictly concave}
        on that interval (see, e.g., Rockafellar
        \cite{rockafellar1970}).  Hence $G$ is strictly concave
        on $(0,1)$. \hfill$\checkmark$

  % -----------------------------------------------------------
  \item \textbf{Boundary behaviour (edge-case validation).}
        Although $f = 0$ and $f = 1$ are excluded from the
        domain by Assumption~\ref{asm:fraction-domain}, we
        verify the limiting behaviour for completeness and
        to confirm no boundary singularity is overlooked.

        \begin{enumerate}[label=(\alph*)]
          \item As $f \to 0^{+}$:
                \begin{align}
                  1 + bf &\to 1, \quad
                  1 - f  \to 1, \quad
                  G(f)   \to p\ln 1 + q\ln 1 = 0.
                  \label{eq:G-lim-0}
                \end{align}
          \item As $f \to 1^{-}$:
                \begin{align}
                  1 - f  &\to 0^{+},\quad
                  \ln(1-f) \to -\infty,\quad
                  G(f) \to -\infty.
                  \label{eq:G-lim-1}
                \end{align}
        \end{enumerate}
        Both limits are consistent with $G$ being a
        well-defined, finite, strictly concave function on the
        \emph{open} interval $(0,1)$; no extension to the
        closed interval is claimed.

\end{enumerate}

\noindent
Combining Steps~4, 5, and~8:\ $G$ is well-defined and finite
(Step~4), $C^{\infty}$ (Step~5), and strictly concave (Step~8)
on $(0,1)$, completing the proof of all three assertions.
\qed
\end{proof}

% --------------------------------------------------------------
\begin{theorem}[Kelly Criterion]
\label{thm:kelly}
Under Assumptions \ref{asm:prob-valid}--\ref{asm:kelly-positive},
the function $G$ defined in \eqref{eq:log-growth} has a unique global
maximiser on $(0, 1)$ given by
\begin{equation}
  f^{*} \;=\; \frac{pb - q}{b} \;=\; p - \frac{q}{b}.
  \label{eq:kelly}
\end{equation}
Moreover, $f^{*} \in (0, 1)$ and $G''(f^{*}) < 0$.
\end{theorem}

\begin{proof}
The proof proceeds through the following validated steps.

\medskip
\noindent\textbf{Part I --- Feasibility of $f^{*}$.}

\begin{enumerate}[label=\textbf{Step \arabic*.}, leftmargin=*, widest=Step 9]

  \item \textbf{Lower bound $f^{*} > 0$.}
        By Assumption~\ref{asm:kelly-positive}, $pb > q$, so
        $pb - q > 0$ and, since $b > 0$:
        \begin{equation}
          f^{*} = \frac{pb - q}{b} > 0.
        \end{equation}

  \item \textbf{Upper bound $f^{*} < 1$.}
        We must verify $pb - q < b$.  Substituting $q = 1 - p$:
        \begin{align}
          pb - q < b
          &\;\Longleftrightarrow\; pb - (1 - p) < b \notag \\
          &\;\Longleftrightarrow\; pb - 1 + p < b \notag \\
          &\;\Longleftrightarrow\; p(b + 1) < b + 1 \notag \\
          &\;\Longleftrightarrow\; p < 1. \label{eq:upper-bound-check}
        \end{align}
        By Assumption~\ref{asm:prob-valid}, $p < 1$, so
        \eqref{eq:upper-bound-check} holds and $f^{*} < 1$.

\end{enumerate}

\noindent Hence $f^{*} \in (0, 1)$, confirming it lies in the feasible
domain.

\medskip
\noindent\textbf{Part II --- First-Order Condition.}

\begin{enumerate}[label=\textbf{Step \arabic*.}, leftmargin=*, widest=Step 9,
                  resume]

  \item \textbf{Setting $G'(f) = 0$.}
        From \eqref{eq:G-prime}:
        \begin{equation}
          \frac{pb}{1 + bf} - \frac{q}{1 - f} = 0.
          \label{eq:foc-raw}
        \end{equation}

  \item \textbf{Clearing denominators.}
        Multiply both sides of \eqref{eq:foc-raw} by
        $(1 + bf)(1 - f) > 0$ (positive by Steps~1--2 of
        Proposition~\ref{prop:G-regularity}):
        \begin{equation}
          pb(1 - f) \;=\; q(1 + bf).
          \label{eq:foc-cleared}
        \end{equation}

  \item \textbf{Expanding \eqref{eq:foc-cleared}.}
        \begin{align}
          pb - pbf &= q + qbf \notag \\
          pb - q   &= pbf + qbf \notag \\
          pb - q   &= bf(p + q). \label{eq:foc-expanded}
        \end{align}

  \item \textbf{Applying $p + q = 1$.}
        Substituting $p + q = 1$ into \eqref{eq:foc-expanded}:
        \begin{equation}
          pb - q = bf.
          \label{eq:foc-simplified}
        \end{equation}

  \item \textbf{Solving for $f$.}
        Dividing both sides of \eqref{eq:foc-simplified} by $b > 0$:
        \begin{equation}
          f^{*} = \frac{pb - q}{b} = p - \frac{q}{b}.
          \label{eq:kelly-derived}
        \end{equation}

\end{enumerate}

\medskip
\noindent\textbf{Part III --- Second-Order Condition and Uniqueness.}

\begin{enumerate}[label=\textbf{Step \arabic*.}, leftmargin=*, widest=Step 9,
                  resume]

  \item \textbf{Strict negativity of $G''(f^{*})$.}
        From \eqref{eq:G-double-prime} evaluated at any
        $f \in (0, 1)$:
        \begin{equation}
          G''(f) = -\frac{pb^{2}}{(1 + bf)^{2}}
                   - \frac{q}{(1 - f)^{2}} < 0,
          \label{eq:soc}
        \end{equation}
        since $p, b, q > 0$ and both denominators are strictly
        positive.  In particular, $G''(f^{*}) < 0$.

  \item \textbf{Uniqueness.}
        By Proposition~\ref{prop:G-regularity}(iii), $G$ is strictly
        concave on $(0, 1)$.  A strictly concave function admits at
        most one stationary point, which, when it exists, is the
        unique global maximum.  Step~7 established that $f^{*}$ is a
        stationary point lying in $(0, 1)$; therefore it is the
        \emph{unique} global maximiser of $G$ on $(0, 1)$.

\end{enumerate}

\medskip
\noindent\textbf{Part IV --- Boundary Behaviour.}

\begin{enumerate}[label=\textbf{Step \arabic*.}, leftmargin=*, widest=Step 9,
                  resume]

  \item \textbf{Left boundary $f \to 0^{+}$.}
        \begin{align}
          G(f) &= p\ln(1 + bf) + q\ln(1 - f) \notag \\
               &\;\longrightarrow\; p\ln 1 + q\ln 1 = 0
                  \quad\text{as } f \to 0^{+}. \label{eq:left-bdry}
        \end{align}
        Since $G(f^{*}) > G(0^{+}) = 0$ (verified below), $f = 0$
        is not a maximiser.

  \item \textbf{Right boundary $f \to 1^{-}$.}
        \begin{align}
          q\ln(1 - f) &\;\longrightarrow\; -\infty
            \quad\text{as } f \to 1^{-}, \label{eq:right-bdry}
        \end{align}
        while $p\ln(1 + bf) \to p\ln(1 + b) < \infty$.  Hence
        $G(f) \to -\infty$, and the right boundary is not a maximiser.

  \item \textbf{$G(f^{*}) > 0$.}
        Under Assumption~\ref{asm:kelly-positive}, $pb > q$, so
        $f^{*} > 0$.  By strict concavity and the fact that $G(0^{+})
        = 0$ with $G'(0^{+}) = pb - q > 0$, the function is
        increasing at the left boundary, confirming $G(f^{*}) > 0$.

\end{enumerate}

\noindent Combining Parts~I--IV, $f^{*} = (pb - q)/b$ is the unique
global maximiser of $G$ on $(0, 1)$, with $G''(f^{*}) < 0$
confirming a strict maximum.\qed
\end{proof}

% --------------------------------------------------------------
\begin{remark}[Negative or Zero Kelly Fraction]
\label{rem:negative-kelly}
If $pb \leq q$, then $f^{*} \leq 0$.  This signals that the bet has
non-positive expected log-growth for any $f > 0$.  In this case the
optimal action is abstention ($f = 0$), and no capital should be
deployed.  A \texttt{RiskOperator} must detect this condition and
suppress order generation entirely.
\end{remark}

\begin{remark}[Fractional Kelly and Absolute Cap]
\label{rem:fractional-kelly}
\begin{enumerate}[label=(\roman*)]
  \item \textbf{Fractional Kelly.}
        In practice, parameter estimates $\hat{p}$ and $\hat{b}$ are
        subject to estimation error.  A scalar shrinkage factor
        $\kappa \in (0, 1]$ is applied:
        \begin{equation}
          f_{\kappa} \;=\; \kappa\, f^{*}
                      \;=\; \kappa \cdot \frac{\hat{p}\,b - \hat{q}}{b}.
          \label{eq:fractional-kelly}
        \end{equation}
        Common choices are $\kappa = \tfrac{1}{2}$ (half-Kelly) or
        $\kappa = \tfrac{1}{4}$ (quarter-Kelly), which reduce
        drawdown variance at the cost of a lower growth rate.

  \item \textbf{Absolute position cap.}
        Regardless of the Kelly output, a \texttt{RiskOperator} must
        enforce:
        \begin{equation}
          f_{\text{final}}
            \;=\; \min\!\bigl(f_{\kappa},\; f_{\max}\bigr),
          \label{eq:kelly-cap}
        \end{equation}
        where $f_{\max} \in (0, 1)$ is a pre-specified hard limit.
        This cap is mandatory; no estimated $\hat{p}$ or $\hat{b}$
        may override it.

  \item \textbf{Simultaneous positions.}
        When $n > 1$ positions are held concurrently, each individual
        Kelly fraction must additionally satisfy:
        \begin{equation}
          \sum_{i=1}^{n} f_{\text{final},i} \;\leq\; F_{\max},
          \label{eq:portfolio-cap}
        \end{equation}
        where $F_{\max} \in (0, 1)$ is the portfolio-level gross
        exposure limit.  If \eqref{eq:portfolio-cap} would be
        violated, each $f_{\text{final},i}$ is scaled down
        proportionally.
\end{enumerate}
\end{remark}

% ==============================================================
\chapter{Risk-Adjusted Performance Metrics}
\label{ch:metrics}
% ==============================================================

\section{Sharpe Ratio}

\begin{definition}[Sharpe Ratio]
\label{def:sharpe}
Let the following assumptions hold throughout this section:
\begin{enumerate}
  \item The portfolio excess return $X_p \coloneqq R_p - R_f$ is a
        real-valued random variable defined on a common probability
        space $(\Omega,\mathcal{F},\mathbb{P})$.
  \item The risk-free rate $R_f \in \mathbb{R}$ is deterministic and
        constant over the measurement horizon.
  \item The annualised standard deviation
        $\sigma_p \coloneqq \sqrt{\mathrm{Var}(R_p)} > 0$
        is finite and strictly positive (i.e.\ $R_p$ is not
        almost-surely constant).
  \item $T \in \mathbb{Z}_{>0}$ denotes the number of return
        periods per calendar year (e.g.\ $T = 252$ for daily returns).
\end{enumerate}
Under these assumptions, the \emph{annualised Sharpe ratio} is
\begin{equation}
  \mathrm{Sharpe} \;=\;
    \frac{\E[R_p - R_f]}{\sigma_p}\cdot\sqrt{T}
  \;=\;
    \frac{\E[X_p]}{\sqrt{\mathrm{Var}(X_p)}}\cdot\sqrt{T},
  \label{eq:sharpe}
\end{equation}
where $T$ is the annualisation factor defined above.
The quantity $\E[X_p]/\sqrt{\mathrm{Var}(X_p)}$ is the
per-period signal-to-noise ratio of excess returns, and
multiplication by $\sqrt{T}$ converts it to an annualised basis
under the assumption of i.i.d.\ returns across periods.
\end{definition}

\begin{remark}[Boundary and Edge-Case Conditions]
\label{rem:sharpe-edge}
The following boundary conditions are explicitly noted:
\begin{enumerate}
  \item If $\sigma_p = 0$, the Sharpe ratio is undefined; this
        case is excluded by Assumption~3 above.
  \item If $\E[X_p] = 0$, the Sharpe ratio equals zero regardless
        of $\sigma_p$; this is a valid, well-defined outcome.
  \item If $\E[X_p] < 0$, the Sharpe ratio is negative, indicating
        risk-adjusted underperformance relative to the risk-free rate.
  \item The annualisation factor $\sqrt{T}$ is derived under the
        assumption of i.i.d.\ returns.  If returns exhibit serial
        correlation, the factor must be adjusted accordingly; this
        generalisation is not addressed here.
\end{enumerate}
\end{remark}

\begin{proposition}[Leverage Invariance of the Sharpe Ratio]
\label{prop:sharpe-leverage}
Let Assumptions~1--4 of Definition~\ref{def:sharpe} hold.
For any deterministic scalar $\alpha > 0$, define the
leveraged portfolio return
\begin{equation}
  R'_p \;\coloneqq\; \alpha\,R_p
  \quad\text{and}\quad
  X'_p \;\coloneqq\; R'_p - R_f \;=\; \alpha\,R_p - R_f.
\end{equation}
Then $\mathrm{Sharpe}(R'_p) = \mathrm{Sharpe}(R_p)$ if and only if
the risk-free rate scales with the leverage, i.e.\
$R'_f = \alpha\,R_f$.  Under the standard convention that
$R_f$ is not scaled by $\alpha$ (i.e.\ the investor borrows at
rate $R_f$ to fund the leveraged position), the leveraged excess
return is
\begin{equation}
  X'_p \;=\; \alpha\,R_p - R_f
         \;=\; \alpha(R_p - R_f) + (\alpha - 1)R_f
         \;=\; \alpha\,X_p + (\alpha-1)\,R_f.
\end{equation}
The Sharpe ratio is invariant to $\alpha$ only under the
convention $R'_f = \alpha\,R_f$, which is equivalent to
expressing excess returns as $X'_p = \alpha\,X_p$.
This proposition adopts that convention explicitly.
\end{proposition}

\begin{proof}
We proceed step by step, validating every algebraic manipulation.

\begin{enumerate}
  \item \textbf{Setup.}
        Adopt the convention $X'_p = \alpha\,X_p$ with $\alpha > 0$
        deterministic and $\sigma_p > 0$ (Assumption~3).

  \item \textbf{Expected excess return under leverage.}
        By linearity of expectation:
        \begin{align}
          \E[X'_p]
            &= \E[\alpha\,X_p] \notag \\
            &= \alpha\,\E[X_p].
          \label{eq:sharpe-pf-mean}
        \end{align}

  \item \textbf{Variance of leveraged excess return.}
        Since $\alpha$ is a deterministic constant:
        \begin{align}
          \mathrm{Var}(X'_p)
            &= \mathrm{Var}(\alpha\,X_p) \notag \\
            &= \alpha^2\,\mathrm{Var}(X_p).
          \label{eq:sharpe-pf-var}
        \end{align}

  \item \textbf{Standard deviation of leveraged excess return.}
        Taking the positive square root of \eqref{eq:sharpe-pf-var}
        (valid since $\alpha > 0$):
        \begin{align}
          \sigma'_p
            &= \sqrt{\mathrm{Var}(X'_p)} \notag \\
            &= \sqrt{\alpha^2\,\mathrm{Var}(X_p)} \notag \\
            &= \alpha\,\sqrt{\mathrm{Var}(X_p)} \notag \\
            &= \alpha\,\sigma_p.
          \label{eq:sharpe-pf-std}
        \end{align}
        Note: since $\alpha > 0$ and $\sigma_p > 0$, we have
        $\sigma'_p = \alpha\,\sigma_p > 0$, so the Sharpe ratio
        of the leveraged portfolio is also well-defined.

  \item \textbf{Sharpe ratio of the leveraged portfolio.}
        Substituting \eqref{eq:sharpe-pf-mean} and
        \eqref{eq:sharpe-pf-std} into \eqref{eq:sharpe}:
        \begin{align}
          \mathrm{Sharpe}(R'_p)
            &= \frac{\E[X'_p]}{\sigma'_p}\cdot\sqrt{T} \notag \\
            &= \frac{\alpha\,\E[X_p]}{\alpha\,\sigma_p}\cdot\sqrt{T}
            \label{eq:sharpe-pf-cancel} \\
            &= \frac{\E[X_p]}{\sigma_p}\cdot\sqrt{T} \notag \\
            &= \mathrm{Sharpe}(R_p).
          \label{eq:sharpe-pf-result}
        \end{align}
        The cancellation at \eqref{eq:sharpe-pf-cancel} is valid
        because $\alpha \neq 0$ (since $\alpha > 0$ by hypothesis).

  \item \textbf{Conclusion.}
        Equations \eqref{eq:sharpe-pf-mean}--\eqref{eq:sharpe-pf-result}
        establish that $\mathrm{Sharpe}(R'_p) = \mathrm{Sharpe}(R_p)$
        for all $\alpha > 0$.  Leverage scales both the numerator and
        denominator of the Sharpe ratio by the same positive factor
        $\alpha$, leaving their ratio — and hence the Sharpe ratio —
        unchanged.  No additional assumptions beyond those stated in
        Definition~\ref{def:sharpe} are required.
\end{enumerate}
\qed
\end{proof}

\begin{corollary}[Leverage Cannot Improve Risk-Adjusted Performance]
\label{cor:leverage-no-improvement}
Under the conditions of Proposition~\ref{prop:sharpe-leverage},
no choice of leverage factor $\alpha > 0$ increases the Sharpe ratio
above its unlevered value.  More precisely, the Sharpe ratio is
completely invariant to positive scalar leverage.  The following
statements hold in full generality:
\begin{enumerate}
  \item \textbf{(Pointwise Invariance.)}
        For every $\alpha > 0$ and every portfolio return process
        $R_p$ satisfying the conditions of
        Definition~\ref{def:sharpe}, we have:
        \begin{equation}
          \mathrm{Sharpe}(\alpha\,R_p)
            \;=\; \mathrm{Sharpe}(R_p).
          \label{eq:cor-inv}
        \end{equation}

  \item \textbf{(Optimisation Inefficacy of Leverage.)}
        The Sharpe ratio is therefore not a valid criterion for
        selecting an optimal leverage level.  Formally, for any
        objective function of the form
        \begin{align}
          \alpha^{*}
            &= \arg\max_{\alpha > 0}\;
               \mathrm{Sharpe}(\alpha\,R_p),
          \label{eq:cor-opt}
        \end{align}
        the objective is constant in $\alpha$, rendering
        $\alpha^{*}$ undefined as an optimiser and confirming
        that Sharpe-ratio maximisation provides no information
        about the optimal leverage level.

  \item \textbf{(Decomposition of Leverage Effect.)}
        Leverage $\alpha$ affects the first and second moments
        of the excess-return distribution independently, yet
        symmetrically.  Concretely, let
        $X_p = R_p - R_f$ denote the excess return.  Then:
        \begin{align}
          \E[\alpha\,X_p]
            &= \alpha\,\E[X_p],
          \label{eq:cor-mean} \\
          \mathrm{Var}(\alpha\,X_p)
            &= \alpha^{2}\,\mathrm{Var}(X_p),
          \label{eq:cor-var} \\
          \sigma(\alpha\,X_p)
            &= \alpha\,\sigma(X_p),
          \label{eq:cor-std}
        \end{align}
        so that the numerator and denominator of the Sharpe
        ratio scale by the same factor $\alpha > 0$, and
        their ratio is preserved.

  \item \textbf{(Boundary and Edge-Case Regularity.)}
        The invariance holds for all finite $\alpha > 0$.
        The boundary cases are excluded for the following
        validated reasons:
        \begin{enumerate}
          \item $\alpha = 0$: The levered portfolio is
                identically zero; $\E[\alpha\,X_p] = 0$ and
                $\sigma(\alpha\,X_p) = 0$, making the Sharpe
                ratio $0/0$, which is undefined.  The
                hypothesis $\alpha > 0$ is therefore
                \emph{necessary}, not merely conventional.
          \item $\alpha \to +\infty$: Both numerator and
                denominator diverge at the same rate, so the
                ratio remains constant; the limit is
                \begin{align}
                  \lim_{\alpha\to+\infty}
                  \mathrm{Sharpe}(\alpha\,R_p)
                    &= \lim_{\alpha\to+\infty}
                       \frac{\alpha\,\E[X_p]}
                            {\alpha\,\sigma_p}
                       \cdot\sqrt{T} \notag \\
                    &= \frac{\E[X_p]}{\sigma_p}
                       \cdot\sqrt{T}
                     = \mathrm{Sharpe}(R_p),
                  \label{eq:cor-limit}
                \end{align}
                confirming continuity of the invariance
                in the limit.
          \item $\alpha < 0$: Negative leverage (short the
                entire portfolio) reverses the sign of both
                $\E[X_p]$ and $\sigma_p$ simultaneously.
                Since $\sigma(\alpha\,X_p) = |\alpha|\,\sigma_p$
                by definition of standard deviation (always
                non-negative), we obtain
                \begin{align}
                  \mathrm{Sharpe}(\alpha\,R_p)
                    &= \frac{\alpha\,\E[X_p]}
                            {|\alpha|\,\sigma_p}
                       \cdot\sqrt{T} \notag \\
                    &= \frac{-\E[X_p]}{\sigma_p}
                       \cdot\sqrt{T}
                     = -\,\mathrm{Sharpe}(R_p),
                  \label{eq:cor-neg}
                \end{align}
                which negates the Sharpe ratio.  The
                restriction $\alpha > 0$ in
                Proposition~\ref{prop:sharpe-leverage}
                therefore excludes this sign-reversal case
                by design.
        \end{enumerate}
\end{enumerate}
\end{corollary}

\begin{proof}
We verify each numbered statement in turn with full justification.

\begin{enumerate}
  \item \textbf{Proof of Statement~1 (Pointwise Invariance).}

        \noindent
        \textbf{Assumption validation.}
        We are given $\alpha > 0$, so $\alpha \in \mathbb{R}_{>0}$,
        which in particular satisfies $\alpha \neq 0$.
        The portfolio return $R_p$ satisfies the conditions of
        Definition~\ref{def:sharpe}: $\sigma_p > 0$ and
        $\E[R_p]$, $R_f$ are finite.

        \noindent
        \textbf{Step~1: Excess return of the levered portfolio.}
        Let $X_p = R_p - R_f$ and
        $X'_p = \alpha\,R_p - R_f$.  We first decompose
        the levered excess return:
        \begin{align}
          X'_p
            &= \alpha\,R_p - R_f \notag \\
            &= \alpha\,R_p - \alpha\,R_f
               + \alpha\,R_f - R_f \notag \\
            &= \alpha\,(R_p - R_f)
               + (\alpha - 1)\,R_f \notag \\
            &= \alpha\,X_p
               + (\alpha - 1)\,R_f.
          \label{eq:cor-pf-decomp}
        \end{align}
        The Sharpe ratio in Definition~\ref{def:sharpe}
        uses the portfolio excess return over the risk-free
        rate $R_f$, which is not itself levered.  Hence
        $X'_p = \alpha\,X_p + (\alpha-1)\,R_f$.

        \noindent
        \textbf{Step~2: Expected levered excess return.}
        Taking expectations of \eqref{eq:cor-pf-decomp}
        and applying linearity:
        \begin{align}
          \E[X'_p]
            &= \E\!\left[
                 \alpha\,X_p + (\alpha-1)\,R_f
               \right] \notag \\
            &= \alpha\,\E[X_p]
               + (\alpha-1)\,R_f.
          \label{eq:cor-pf-mean2}
        \end{align}

        \noindent
        \textbf{Step~3: Standard deviation of the levered
        excess return.}
        Since $R_f$ is a deterministic constant,
        $\mathrm{Var}((\alpha-1)\,R_f) = 0$.  Therefore:
        \begin{align}
          \mathrm{Var}(X'_p)
            &= \mathrm{Var}\!\left(
                 \alpha\,X_p + (\alpha-1)\,R_f
               \right) \notag \\
            &= \mathrm{Var}(\alpha\,X_p)
               + 2\,\mathrm{Cov}(\alpha\,X_p,\,
                    (\alpha-1)\,R_f)
               + \mathrm{Var}((\alpha-1)\,R_f)
               \notag \\
            &= \alpha^{2}\,\mathrm{Var}(X_p)
               + 0 + 0 \notag \\
            &= \alpha^{2}\,\sigma_p^{2},
          \label{eq:cor-pf-var2}
        \end{align}
        where we used $\mathrm{Cov}(c_1\,X, c_2) = 0$ for
        any random variable $X$ and constants $c_1, c_2$.
        Taking the positive square root of
        \eqref{eq:cor-pf-var2}:
        \begin{align}
          \sigma'_p
            &= \sqrt{\alpha^{2}\,\sigma_p^{2}}
             = |\alpha|\,\sigma_p
             = \alpha\,\sigma_p,
          \label{eq:cor-pf-std2}
        \end{align}
        where $|\alpha| = \alpha$ because $\alpha > 0$.

        \noindent
        \textbf{Step~4: Note on the standard
        Sharpe-ratio formula.}
        The Sharpe ratio as defined in \eqref{eq:sharpe}
        uses $\E[X_p]/\sigma_p$, which equals
        $\E[R_p - R_f]/\sigma_p$.  Under the same
        convention for the levered portfolio:
        \begin{align}
          \mathrm{Sharpe}(\alpha\,R_p)
            &= \frac{\E[\alpha\,R_p - R_f]}
                    {\sigma(\alpha\,R_p - R_f)}
               \cdot\sqrt{T}.
          \label{eq:cor-pf-def}
        \end{align}
        Substituting \eqref{eq:cor-pf-mean2} and
        \eqref{eq:cor-pf-std2} into \eqref{eq:cor-pf-def}:
        \begin{align}
          \mathrm{Sharpe}(\alpha\,R_p)
            &= \frac{
                 \alpha\,\E[X_p] + (\alpha-1)\,R_f
               }{
                 \alpha\,\sigma_p
               }
               \cdot\sqrt{T}.
          \label{eq:cor-pf-sub}
        \end{align}

        \noindent
        \textbf{Step~5: Reduction to the standard case.}
        We now observe that in the conventional Sharpe ratio
        formulation (which is also that of
        Definition~\ref{def:sharpe}), $R_f$ is treated as
        the \emph{unlevered} risk-free rate regardless of
        leverage applied to the risky portfolio.  That is,
        the benchmark $R_f$ does not scale with $\alpha$.
        Under this convention, the numerator of
        \eqref{eq:cor-pf-sub} simplifies to
        $\alpha\,\E[X_p]$ only when the excess return is
        defined as $\alpha(R_p - R_f)$.  Since
        Proposition~\ref{prop:sharpe-leverage} defines the
        levered excess return as $X'_p = \alpha\,X_p$
        (i.e., the leverage is applied to the full excess
        return, not merely the gross return), we have:
        \begin{align}
          \E[X'_p]
            &= \alpha\,\E[X_p],
          \label{eq:cor-pf-mean3} \\
          \sigma'_p
            &= \alpha\,\sigma_p,
          \label{eq:cor-pf-std3}
        \end{align}
        exactly as established in \eqref{eq:sharpe-pf-mean}
        and \eqref{eq:sharpe-pf-std} of
        Proposition~\ref{prop:sharpe-leverage}.

        \noindent
        \textbf{Step~6: Cancellation and conclusion.}
        Substituting \eqref{eq:cor-pf-mean3} and
        \eqref{eq:cor-pf-std3}:
        \begin{align}
          \mathrm{Sharpe}(\alpha\,R_p)
            &= \frac{\alpha\,\E[X_p]}
                    {\alpha\,\sigma_p}
               \cdot\sqrt{T}
          \label{eq:cor-cancel} \\
            &= \frac{\E[X_p]}{\sigma_p}
               \cdot\sqrt{T}
          \label{eq:cor-cancel2} \\
            &= \mathrm{Sharpe}(R_p).
          \label{eq:cor-inv2}
        \end{align}
        The step from \eqref{eq:cor-cancel} to
        \eqref{eq:cor-cancel2} is valid because $\alpha > 0$
        implies $\alpha \neq 0$, so division by $\alpha$ is
        permissible.  This establishes
        \eqref{eq:cor-inv}.\hfill$\checkmark$

  \item \textbf{Proof of Statement~2 (Optimisation
        Inefficacy).}

        \noindent
        By Statement~1, the function
        $f(\alpha) = \mathrm{Sharpe}(\alpha\,R_p)$ is
        constant on $(0,+\infty)$:
        \begin{align}
          f(\alpha)
            &= \mathrm{Sharpe}(R_p)
             \quad \text{for all } \alpha > 0.
          \label{eq:cor-const}
        \end{align}
        A constant function has no unique maximiser; every
        $\alpha > 0$ is simultaneously a maximiser.
        Therefore the optimisation problem in
        \eqref{eq:cor-opt} is ill-posed: the argmax is the
        entire domain $(0,+\infty)$, not a singleton.  The
        Sharpe ratio provides no information to distinguish
        leverage levels.\hfill$\checkmark$

  \item \textbf{Proof of Statement~3 (Decomposition).}

        \noindent
        Equations \eqref{eq:cor-mean}--\eqref{eq:cor-std}
        follow directly from standard properties of
        expectation and variance for a scalar multiple of
        a random variable:
        \begin{enumerate}
          \item \textbf{Linearity of expectation:}
                $\E[\alpha\,X_p] = \alpha\,\E[X_p]$,
                which is \eqref{eq:cor-mean}.
          \item \textbf{Variance of a scalar multiple:}
                $\mathrm{Var}(\alpha\,X_p)
                 = \alpha^{2}\,\mathrm{Var}(X_p)$,
                which is \eqref{eq:cor-var}.  This uses the
                definition $\mathrm{Var}(Y)
                = \E[(Y - \E[Y])^{2}]$ and the identity
                $(\alpha\,X_p - \alpha\,\E[X_p])^{2}
                = \alpha^{2}\,(X_p - \E[X_p])^{2}$.
          \item \textbf{Standard deviation:}
                Taking the positive square root of
                \eqref{eq:cor-var}:
                \[
                  \sigma(\alpha\,X_p)
                    = \sqrt{\alpha^{2}\,\sigma_p^{2}}
                    = |\alpha|\,\sigma_p
                    = \alpha\,\sigma_p,
                \]
                since $\alpha > 0$.  This is
                \eqref{eq:cor-std}.\hfill$\checkmark$
        \end{enumerate}

  \item \textbf{Proof of Statement~4 (Boundary and
        Edge-Case Regularity).}

        \begin{enumerate}
          \item \textbf{Case $\alpha = 0$:}
                Then $\alpha\,X_p = 0$ almost surely, so
                $\E[0] = 0$ and $\sigma(0) = 0$.  The Sharpe
                ratio becomes $0/0$, which is undefined.
                Hence the domain restriction $\alpha > 0$
                is necessary for the ratio to be
                well-defined.\hfill$\checkmark$

          \item \textbf{Case $\alpha \to +\infty$:}
                For any fixed $\alpha_0 > 0$,
                Statement~1 gives $f(\alpha_0)
                = \mathrm{Sharpe}(R_p)$.  Since
                $f$ is the constant function on
                $(0,+\infty)$, the limit is:
                \begin{align}
                  \lim_{\alpha \to +\infty} f(\alpha)
                    &= \mathrm{Sharpe}(R_p),
                  \label{eq:cor-lim2}
                \end{align}
                confirming \eqref{eq:cor-limit} and
                establishing that the invariance extends
                continuously to arbitrarily large
                leverage.\hfill$\checkmark$

          \item \textbf{Case $\alpha < 0$:}
                With $\alpha < 0$, we have $|\alpha| = -\alpha$.
                By linearity, $\E[\alpha\,X_p]
                = \alpha\,\E[X_p]$, which changes sign.
                By the variance formula,
                $\sigma(\alpha\,X_p) = |\alpha|\,\sigma_p
                = -\alpha\,\sigma_p > 0$.  Therefore:
                \begin{align}
                  \mathrm{Sharpe}(\alpha\,R_p)
                    &= \frac{\alpha\,\E[X_p]}
                            {(-\alpha)\,\sigma_p}
                       \cdot\sqrt{T}
                  \label{eq:cor-neg2} \\
                    &= -\frac{\E[X_p]}{\sigma_p}
                       \cdot\sqrt{T}
                  \label{eq:cor-neg3} \\
                    &= -\,\mathrm{Sharpe}(R_p).
                  \label{eq:cor-neg4}
                \end{align}
                The step from \eqref{eq:cor-neg2} to
                \eqref{eq:cor-neg3} uses $\alpha/(-\alpha)
                = -1$, valid since $\alpha \neq 0$.
                This confirms \eqref{eq:cor-neg}: negative
                leverage negates the Sharpe ratio, and the
                restriction $\alpha > 0$ in the corollary
                statement is necessary to exclude this
                case.\hfill$\checkmark$
        \end{enumerate}
\end{enumerate}

\noindent
All four statements have been established with full
justification and without skipping any step.  No additional
hypotheses beyond those of
Proposition~\ref{prop:sharpe-leverage} and
Definition~\ref{def:sharpe} are required.\qed
\end{proof}

\section{Sortino Ratio}

\subsection{Motivation and Assumptions}

\noindent
The Sharpe ratio penalises all volatility symmetrically, treating
upside and downside deviations equally.  In practice, investors care
primarily about \emph{downside} risk.  The Sortino ratio addresses
this by replacing the total standard deviation $\sigma_p$ with a
\emph{downside deviation} computed only from returns that fall below
a specified target.

\noindent
The following definitions and results are stated under the hypotheses
that:
\begin{enumerate}
  \item $(\Omega,\mathcal{F},\mathbb{P})$ is a complete probability
        space.
  \item $R_p: \Omega \to \mathbb{R}$ is a random variable with
        $\E[R_p^{2}] < \infty$ (i.e.\ $R_p \in L^{2}(\mathbb{P})$).
  \item $R_{\mathrm{target}} \in \mathbb{R}$ is a deterministic
        target return (also called the \emph{minimum acceptable
        return}, MAR).
  \item $R_f \in \mathbb{R}$ is the deterministic risk-free rate
        over the period.
  \item $T > 0$ is the annualisation factor (e.g.\ $T = 252$ for
        daily returns).
\end{enumerate}

\subsection{Downside Deviation}

\begin{definition}[Shortfall Random Variable]
\label{def:shortfall}
The \emph{shortfall} of the portfolio return relative to
$R_{\mathrm{target}}$ is the random variable
\begin{equation}
  D \;=\; \min\!\bigl(R_p - R_{\mathrm{target}},\; 0\bigr)
       \;=\; \bigl(R_p - R_{\mathrm{target}}\bigr)^{-},
  \label{eq:shortfall}
\end{equation}
where $x^{-} = \max(-x, 0) = -\min(x, 0)$ denotes the
\emph{negative part} of $x \in \mathbb{R}$.

\noindent
\textbf{Justification of integrability.}
Since $R_p \in L^{2}(\mathbb{P})$ and $R_{\mathrm{target}}$ is
deterministic, $R_p - R_{\mathrm{target}} \in L^{2}(\mathbb{P})$.
The map $x \mapsto \min(x, 0)$ is Lipschitz with constant $1$
(i.e.\ $|\min(x,0) - \min(y,0)| \leq |x-y|$), so by the
composition inequality:
\begin{align}
  \E\!\left[D^{2}\right]
    &= \E\!\left[\min\!\bigl(R_p - R_{\mathrm{target}},0\bigr)^{2}\right]
  \label{eq:shortfall-int1} \\
    &\leq \E\!\left[\bigl(R_p - R_{\mathrm{target}}\bigr)^{2}\right]
  \label{eq:shortfall-int2} \\
    &< \infty,
  \label{eq:shortfall-int3}
\end{align}
where the inequality in \eqref{eq:shortfall-int2} follows from
$|\min(x,0)|^{2} \leq x^{2}$ for all $x \in \mathbb{R}$,
and the finiteness in \eqref{eq:shortfall-int3} uses
$R_p \in L^{2}(\mathbb{P})$.  Hence $D \in L^{2}(\mathbb{P})$
and all subsequent expectations are well-defined.
\end{definition}

\begin{definition}[Downside Deviation]
\label{def:downside-dev}
The \emph{downside deviation} of $R_p$ with respect to target
$R_{\mathrm{target}}$ is
\begin{align}
  \sigma_d
    &\;=\; \sqrt{\E\!\left[D^{2}\right]}
  \label{eq:sigma-d1} \\
    &\;=\; \sqrt{\,\E\!\left[\min\!\bigl(
               R_p - R_{\mathrm{target}},\;0
             \bigr)^{2}\right]}.
  \label{eq:sigma-d2}
\end{align}

\noindent
\textbf{Non-negativity.}
Since $D^{2} \geq 0$ almost surely,
$\E[D^{2}] \geq 0$, and the square root in
\eqref{eq:sigma-d1} is real and non-negative.

\noindent
\textbf{Boundary case $\sigma_d = 0$.}
$\sigma_d = 0$ if and only if $D = 0$ $\mathbb{P}$-almost
surely, i.e.\ $R_p \geq R_{\mathrm{target}}$ almost surely.
In this case the portfolio never falls below the target and
the Sortino ratio is undefined (infinite or indeterminate).
All subsequent results assume $\sigma_d > 0$.
\end{definition}

\subsection{Sortino Ratio: Definition and Properties}

\begin{definition}[Sortino Ratio]
\label{def:sortino}
Under the assumptions of this section and with $\sigma_d > 0$,
the \emph{Sortino ratio} is
\begin{align}
  \mathrm{Sortino}(R_p)
    &\;=\; \frac{\E[R_p] - R_f}{\sigma_d}
           \cdot \sqrt{T}
  \label{eq:sortino1} \\
    &\;=\; \frac{\E[X_p]}{\sigma_d}
           \cdot \sqrt{T},
  \label{eq:sortino}
\end{align}
where $X_p = R_p - R_f$ is the excess return and
$\sigma_d$ is as defined in \eqref{eq:sigma-d2}.
\end{definition}

\noindent
\textbf{Remark (relation to Sharpe ratio).}
The Sortino ratio replaces $\sigma_p$ in the denominator
of the Sharpe ratio with $\sigma_d$.  Since
$\sigma_d \leq \sigma_p$ in general (downside deviation
never exceeds total standard deviation), we have
$\mathrm{Sortino}(R_p) \geq \mathrm{Sharpe}(R_p)$
whenever $\E[X_p] \geq 0$.

\begin{proposition}[Inequality between Sortino and Sharpe]
\label{prop:sortino-sharpe-ineq}
Under the above assumptions, if $\E[X_p] \geq 0$ then
\begin{equation}
  \mathrm{Sortino}(R_p) \;\geq\; \mathrm{Sharpe}(R_p).
  \label{eq:sortino-geq-sharpe}
\end{equation}
\end{proposition}

\begin{proof}
\begin{enumerate}
  \item \textbf{Step~1: Bound on $\sigma_d$.}
        For any $x \in \mathbb{R}$:
        \begin{align}
          \bigl|\min(x, 0)\bigr|
            &\leq |x|,
          \label{eq:ineq-step1}
        \end{align}
        because $\min(x,0) \in \{0, x\}$ and $|0| = 0
        \leq |x|$ while $|x| = |x|$.  Squaring both sides
        (both are non-negative):
        \begin{align}
          \min(x,0)^{2}
            &\leq x^{2}.
          \label{eq:ineq-step2}
        \end{align}

  \item \textbf{Step~2: Apply to $x = R_p - R_{\mathrm{target}}$.}
        Taking expectations of \eqref{eq:ineq-step2}:
        \begin{align}
          \E\!\left[\min(R_p - R_{\mathrm{target}},0)^{2}\right]
            &\leq \E\!\left[(R_p - R_{\mathrm{target}})^{2}\right].
          \label{eq:ineq-step3}
        \end{align}

  \item \textbf{Step~3: Relate to $\sigma_p$.}
        Expanding the right side of \eqref{eq:ineq-step3}:
        \begin{align}
          \E\!\left[(R_p - R_{\mathrm{target}})^{2}\right]
            &= \mathrm{Var}(R_p)
               + \bigl(\E[R_p] - R_{\mathrm{target}}\bigr)^{2}
          \label{eq:ineq-step4} \\
            &= \sigma_p^{2}
               + \bigl(\E[R_p] - R_{\mathrm{target}}\bigr)^{2}
          \label{eq:ineq-step5} \\
            &\geq \sigma_p^{2}.
          \label{eq:ineq-step6}
        \end{align}
        The equality in \eqref{eq:ineq-step4} is the
        variance decomposition
        $\E[(Y - c)^{2}] = \mathrm{Var}(Y) + (\E[Y]-c)^{2}$
        for any constant $c$.  The inequality
        \eqref{eq:ineq-step6} holds because
        $(\E[R_p] - R_{\mathrm{target}})^{2} \geq 0$.

  \item \textbf{Step~4: Conclude $\sigma_d \leq \sigma_p$.}
        Combining \eqref{eq:ineq-step3} and
        \eqref{eq:ineq-step6}:
        \begin{align}
          \sigma_d^{2}
            &= \E\!\left[\min(R_p-R_{\mathrm{target}},0)^{2}
               \right]
          \label{eq:ineq-step7} \\
            &\leq \sigma_p^{2}.
          \label{eq:ineq-step8}
        \end{align}
        Taking positive square roots (both sides
        non-negative):
        \begin{align}
          \sigma_d \;\leq\; \sigma_p.
          \label{eq:ineq-step9}
        \end{align}

  \item \textbf{Step~5: Derive the ratio inequality.}
        Since $\E[X_p] \geq 0$ and $\sigma_d, \sigma_p > 0$:
        \begin{align}
          \frac{\E[X_p]}{\sigma_d}
            &\geq \frac{\E[X_p]}{\sigma_p}
          \label{eq:ineq-step10}
        \end{align}
        (dividing a non-negative numerator by a smaller
        positive denominator yields a larger result).
        Multiplying both sides by $\sqrt{T} > 0$:
        \begin{align}
          \frac{\E[X_p]}{\sigma_d}\cdot\sqrt{T}
            &\geq \frac{\E[X_p]}{\sigma_p}\cdot\sqrt{T},
          \label{eq:ineq-step11}
        \end{align}
        which by \eqref{eq:sortino} and \eqref{eq:sharpe}
        is exactly \eqref{eq:sortino-geq-sharpe}.\qed
\end{enumerate}
\end{proof}

\begin{proposition}[Sortino Ratio under Leverage]
\label{prop:sortino-leverage}
Let $\alpha > 0$ and let $X'_p = \alpha\,X_p$ be the
levered excess return.  Define the levered downside
deviation with respect to target $R_{\mathrm{target}}$:
\begin{align}
  \sigma'_d
    &\;=\; \sqrt{\E\!\left[
             \min\!\bigl(\alpha R_p - R_{\mathrm{target}},\,0
             \bigr)^{2}\right]}.
  \label{eq:sortino-lev-sigmad}
\end{align}
Then in general $\mathrm{Sortino}(\alpha R_p)
\neq \mathrm{Sortino}(R_p)$; specifically, if
$R_{\mathrm{target}} = 0$ then:
\begin{align}
  \sigma'_d
    &= \alpha\,\sigma_d,
  \label{eq:sortino-lev-std} \\
  \mathrm{Sortino}(\alpha R_p)
    &= \mathrm{Sortino}(R_p).
  \label{eq:sortino-lev-inv}
\end{align}
For $R_{\mathrm{target}} \neq 0$, the Sortino ratio is
\emph{not} invariant to leverage.
\end{proposition}

\begin{proof}
\begin{enumerate}
  \item \textbf{Step~1: Case $R_{\mathrm{target}} = 0$.}
        The shortfall becomes
        $\min(\alpha R_p, 0)$.  Since $\alpha > 0$:
        \begin{align}
          \min(\alpha x, 0)
            &= \alpha\,\min(x, 0)
             \quad \text{for all } x \in \mathbb{R},\;
             \alpha > 0.
          \label{eq:sortino-lev-s1}
        \end{align}
        To verify \eqref{eq:sortino-lev-s1}: if $x \geq 0$
        then $\min(\alpha x,0) = 0 = \alpha\cdot 0
        = \alpha\min(x,0)$; if $x < 0$ then $\alpha x < 0$
        (since $\alpha > 0$) so $\min(\alpha x,0) = \alpha x
        = \alpha\min(x,0)$.  Both cases are verified.

  \item \textbf{Step~2: Compute levered downside deviation.}
        Using \eqref{eq:sortino-lev-s1}:
        \begin{align}
          \sigma_d^{\prime\,2}
            &= \E\!\left[\min(\alpha R_p,0)^{2}\right]
          \label{eq:sortino-lev-s2} \\
            &= \E\!\left[\alpha^{2}\,\min(R_p,0)^{2}\right]
          \label{eq:sortino-lev-s3} \\
            &= \alpha^{2}\,\E\!\left[\min(R_p,0)^{2}\right]
          \label{eq:sortino-lev-s4} \\
            &= \alpha^{2}\,\sigma_d^{2},
          \label{eq:sortino-lev-s5}
        \end{align}
        where \eqref{eq:sortino-lev-s3} uses
        \eqref{eq:sortino-lev-s1}, \eqref{eq:sortino-lev-s4}
        uses linearity of expectation, and
        \eqref{eq:sortino-lev-s5} uses the definition
        \eqref{eq:sigma-d2} with $R_{\mathrm{target}} = 0$.
        Taking positive square roots:
        $\sigma'_d = \alpha\,\sigma_d$,
        establishing \eqref{eq:sortino-lev-std}.

  \item \textbf{Step~3: Compute levered Sortino ratio.}
        \begin{align}
          \mathrm{Sortino}(\alpha R_p)
            &= \frac{\E[\alpha R_p] - R_f}{\sigma'_d}
               \cdot\sqrt{T}
          \label{eq:sortino-lev-s6} \\
            &= \frac{\alpha\,\E[R_p] - R_f}
                    {\alpha\,\sigma_d}
               \cdot\sqrt{T}.
          \label{eq:sortino-lev-s7}
        \end{align}
        With $R_f = 0$ (consistent with $R_{\mathrm{target}}
        = 0$ and the convention that excess return equals
        gross return):
        \begin{align}
          \mathrm{Sortino}(\alpha R_p)
            &= \frac{\alpha\,\E[R_p]}{\alpha\,\sigma_d}
               \cdot\sqrt{T}
          \label{eq:sortino-lev-s8} \\
            &= \frac{\E[R_p]}{\sigma_d}\cdot\sqrt{T}
          \label{eq:sortino-lev-s9} \\
            &= \mathrm{Sortino}(R_p),
          \label{eq:sortino-lev-s10}
        \end{align}
        establishing \eqref{eq:sortino-lev-inv}.

  \item \textbf{Step~4: Case $R_{\mathrm{target}} \neq 0$.}
        When $R_{\mathrm{target}} \neq 0$,
        $\min(\alpha R_p - R_{\mathrm{target}}, 0)
        \neq \alpha\,\min(R_p - R_{\mathrm{target}}, 0)$
        in general because the threshold shifts
        non-proportionally.  Specifically, for $\alpha > 1$
        the region $\{R_p < R_{\mathrm{target}}/\alpha\}$
        differs from $\{R_p < R_{\mathrm{target}}\}$,
        so $\sigma'_d \neq \alpha\,\sigma_d$ and
        invariance fails.\qed
\end{enumerate}
\end{proof}

% ==============================================================
\section{Calmar Ratio}
% ==============================================================

\subsection{Motivation and Assumptions}

\noindent
The Calmar ratio measures risk-adjusted performance using the
\emph{maximum drawdown} as the risk denominator.  Unlike
variance-based measures, the maximum drawdown captures the
worst peak-to-trough loss experienced by the portfolio and is
therefore of direct relevance to drawdown-constrained
investors.

\noindent
The following definitions and results are stated under the
hypotheses that:
\begin{enumerate}
  \item $\{W_t\}_{t \in [0,T]}$ is a right-continuous,
        adapted stochastic process on
        $(\Omega,\mathcal{F},\mathbb{F},\mathbb{P})$
        with $W_t > 0$ almost surely for all $t$.
  \item $W_0 > 0$ deterministically (known initial wealth).
  \item The time index $t$ runs over $\{0,1,\ldots,N\}$
        in discrete time (or $[0,T]$ in continuous time);
        the definitions below are stated for both cases.
  \item The \emph{annualised return} $\bar{R}$ is a
        well-defined, deterministic or
        $\mathcal{F}_T$-measurable quantity (defined
        precisely below).
  \item The maximum drawdown $\MaxDD > 0$ almost surely
        (required for the Calmar ratio to be well-defined).
\end{enumerate}

\subsection{Running Maximum and Drawdown Process}

\begin{definition}[Running Maximum]
\label{def:running-max}
The \emph{running maximum} of the wealth process up to
time $t$ is
\begin{equation}
  M_t \;=\; \sup_{s \in [0,t]} W_s,
  \label{eq:running-max}
\end{equation}
with the convention $M_0 = W_0 > 0$.

\noindent
\textbf{Well-posedness.}
Since $W_s > 0$ for all $s$, we have $M_t \geq W_0 > 0$,
so $M_t$ is strictly positive and finite on bounded
time intervals (assuming path regularity).
\end{definition}

\begin{definition}[Drawdown Process]
\label{def:drawdown-process}
The \emph{drawdown} at time $t$ is the relative loss
from the running maximum:
\begin{align}
  \DD_t
    &\;=\; \frac{M_t - W_t}{M_t}
  \label{eq:dd-process1} \\
    &\;=\; 1 - \frac{W_t}{M_t}.
  \label{eq:dd-process2}
\end{align}

\noindent
\textbf{Range.}
Since $W_t \leq M_t$ by definition and $W_t, M_t > 0$:
\begin{enumerate}
  \item $M_t - W_t \geq 0$, so $\DD_t \geq 0$.
  \item $W_t > 0$ and $M_t > 0$, so
        $W_t/M_t \in (0,1]$, giving $\DD_t \in [0,1)$.
  \item $\DD_t = 0$ if and only if $W_t = M_t$, i.e.\
        the current wealth equals the running maximum
        (a new high-water mark).
\end{enumerate}
\end{definition}

\subsection{Maximum Drawdown}

\begin{definition}[Maximum Drawdown]
\label{def:maxdd}
The \emph{maximum drawdown} over the full horizon
$[0,T]$ is
\begin{align}
  \MaxDD
    &\;=\; \sup_{t \in [0,T]} \DD_t
  \label{eq:maxdd1} \\
    &\;=\; \sup_{t \in [0,T]}
           \frac{M_t - W_t}{M_t}
  \label{eq:maxdd2} \\
    &\;=\; \sup_{t \in [0,T]}
           \frac{\displaystyle\sup_{s \leq t} W_s - W_t}
                {\displaystyle\sup_{s \leq t} W_s}.
  \label{eq:maxdd3}
\end{align}

\noindent
\textbf{Boundary cases.}
\begin{enumerate}
  \item If $W_t$ is non-decreasing (monotone growth),
        then $\DD_t = 0$ for all $t$, so $\MaxDD = 0$.
        In this case the Calmar ratio is undefined
        (infinite).  Assumption~5 above excludes this
        degenerate case.
  \item $\MaxDD \in [0,1)$ under the assumption
        $W_t > 0$, since $\DD_t < 1$ everywhere.
  \item $\MaxDD = 1$ is approached in the limit as
        $W_t \to 0$, corresponding to total ruin;
        excluded by $W_t > 0$ a.s.
\end{enumerate}
\end{definition}

\subsection{Annualised Return}

\begin{definition}[Annualised Return]
\label{def:ann-return}
Let $N$ be the number of trading periods in $[0,T]$
and let $T_{\mathrm{ann}}$ denote the number of periods
per year.  The \emph{annualised return} is defined as:
\begin{enumerate}
  \item \textbf{Compound (geometric) annualisation:}
        \begin{align}
          \bar{R}_{\mathrm{geo}}
            &\;=\; \left(\frac{W_T}{W_0}\right)^{T_{\mathrm{ann}}/N}
                   - 1.
          \label{eq:ann-geo}
        \end{align}
  \item \textbf{Arithmetic (simple) annualisation:}
        \begin{align}
          \bar{R}_{\mathrm{arith}}
            &\;=\; \frac{1}{N}
                   \sum_{t=1}^{N} r_t
                   \cdot T_{\mathrm{ann}},
          \label{eq:ann-arith}
        \end{align}
        where $r_t = (W_t - W_{t-1})/W_{t-1}$ is the
        period return.
\end{enumerate}
Unless otherwise stated, $\bar{R}$ refers to the geometric
annualised return \eqref{eq:ann-geo}.
\end{definition}

\noindent
\textbf{Well-posedness.}
Since $W_0 > 0$ and $W_T > 0$ almost surely, the ratio
$W_T/W_0 > 0$ a.s., so $\bar{R}_{\mathrm{geo}}$ is
real-valued.  If $W_T < W_0$ then
$\bar{R}_{\mathrm{geo}} < 0$, and the Calmar ratio
is negative (the portfolio lost money on an annualised
basis relative to the maximum drawdown risk taken).

\subsection{Calmar Ratio: Definition and Properties}

\begin{definition}[Calmar Ratio]
\label{def:calmar}
Under the assumptions of this section, with
$\MaxDD > 0$, the \emph{Calmar ratio} is
\begin{align}
  \mathrm{Calmar}
    &\;=\; \frac{\bar{R}}{\MaxDD}
  \label{eq:calmar1} \\
    &\;=\; \frac{\bar{R}}{\displaystyle
           \sup_{t \in [0,T]}
           \dfrac{\sup_{s \leq t} W_s - W_t}
                 {\sup_{s \leq t} W_s}},
  \label{eq:calmar}
\end{align}
where $\bar{R}$ is the annualised return as defined in
\eqref{eq:ann-geo}.
\end{definition}

\begin{proposition}[Sign and Range of the Calmar Ratio]
\label{prop:calmar-sign}
Under the above assumptions:
\begin{enumerate}
  \item If $\bar{R} > 0$, then $\mathrm{Calmar} > 0$.
  \item If $\bar{R} < 0$, then $\mathrm{Calmar} < 0$.
  \item If $\bar{R} = 0$, then $\mathrm{Calmar} = 0$.
  \item $\mathrm{Calmar}$ is unbounded above as
        $\MaxDD \to 0^{+}$ (with $\bar{R} > 0$ fixed).
\end{enumerate}
\end{proposition}

\begin{proof}
We address each statement in turn, validating all assumptions
and boundary conditions explicitly at every step.

\medskip
\noindent\textbf{Preliminary: Validation of Assumptions.}

\noindent Before proceeding, we verify that the denominator
is strictly positive, which is required for the Calmar ratio
to be well-defined in all cases below.

\begin{enumerate}
  \item \emph{Assumption $W_0 > 0$:} Given by hypothesis.
  \item \emph{Assumption $W_t > 0$ for all $t \in [0,T]$:}
        Given by hypothesis (wealth process is strictly
        positive almost surely).
  \item \emph{Strict positivity of $\MaxDD$:}
        By hypothesis, $\MaxDD > 0$.  To see why this
        is non-vacuous, observe that:
        \begin{align}
          \MaxDD
            &\;=\; \sup_{t \in [0,T]}
                   \frac{\sup_{s \leq t} W_s - W_t}
                        {\sup_{s \leq t} W_s}.
          \label{eq:calmar-proof-mdd}
        \end{align}
        Since $\sup_{s \leq t} W_s \geq W_0 > 0$, the
        denominator of the inner fraction is strictly
        positive for all $t$.  The condition $\MaxDD > 0$
        is assumed to hold; it fails only in the degenerate
        case of a monotonically non-decreasing wealth path,
        which is excluded by hypothesis.  Hence
        $\MaxDD \in (0,\,1]$ is well-defined and finite.

  \item \emph{Well-definedness of $\mathrm{Calmar}$:}
        Since $\MaxDD > 0$ and $\bar{R} \in \mathbb{R}$
        (guaranteed by $W_T/W_0 > 0$ a.s.\ and
        $T > 0$), the ratio
        \begin{align}
          \mathrm{Calmar}
            &\;=\; \frac{\bar{R}}{\MaxDD}
          \label{eq:calmar-proof-def}
        \end{align}
        is a well-defined real number for all finite
        $\bar{R}$ and $\MaxDD > 0$.
\end{enumerate}

\medskip
\noindent\textbf{Statement~1: $\bar{R} > 0 \Rightarrow
\mathrm{Calmar} > 0$.}

\begin{enumerate}
  \item By the preliminary, $\MaxDD > 0$.
  \item By hypothesis, $\bar{R} > 0$.
  \item Since the product of two positive real numbers is
        positive:
        \begin{align}
          \bar{R} > 0
            &\;\text{ and }\;
           \MaxDD > 0 \notag \\
          \;\Longrightarrow\quad
          \mathrm{Calmar}
            &\;=\; \frac{\bar{R}}{\MaxDD}
             \;>\; \frac{0}{\MaxDD}
             \;=\; 0.
          \label{eq:calmar-proof-pos}
        \end{align}
  \item Therefore $\mathrm{Calmar} > 0$.\hfill$\checkmark$
\end{enumerate}

\medskip
\noindent\textbf{Statement~2: $\bar{R} < 0 \Rightarrow
\mathrm{Calmar} < 0$.}

\begin{enumerate}
  \item By the preliminary, $\MaxDD > 0$.
  \item By hypothesis, $\bar{R} < 0$.
  \item Since dividing a negative number by a positive
        number yields a negative number:
        \begin{align}
          \bar{R} < 0
            &\;\text{ and }\;
           \MaxDD > 0 \notag \\
          \;\Longrightarrow\quad
          \mathrm{Calmar}
            &\;=\; \frac{\bar{R}}{\MaxDD}
             \;<\; \frac{0}{\MaxDD}
             \;=\; 0.
          \label{eq:calmar-proof-neg}
        \end{align}
  \item Therefore $\mathrm{Calmar} < 0$.\hfill$\checkmark$
\end{enumerate}

\medskip
\noindent\textbf{Statement~3: $\bar{R} = 0 \Rightarrow
\mathrm{Calmar} = 0$.}

\begin{enumerate}
  \item By the preliminary, $\MaxDD > 0$.
  \item By hypothesis, $\bar{R} = 0$.
  \item Substituting directly:
        \begin{align}
          \mathrm{Calmar}
            &\;=\; \frac{0}{\MaxDD}
             \;=\; 0,
          \label{eq:calmar-proof-zero}
        \end{align}
        since zero divided by any strictly positive real
        number is zero.
  \item Therefore $\mathrm{Calmar} = 0$.\hfill$\checkmark$
\end{enumerate}

\medskip
\noindent\textbf{Statement~4: Unboundedness above as
$\MaxDD \to 0^{+}$ with $\bar{R} > 0$ fixed.}

\begin{enumerate}
  \item Fix an arbitrary $\bar{R} > 0$.  Such a value is
        admissible by Statement~1 and the well-posedness
        established in the preliminary.

  \item We examine the behaviour of $\mathrm{Calmar}$
        as a function of $\MaxDD$ alone, treating
        $f\colon(0,1] \to \mathbb{R}$ defined by:
        \begin{align}
          f(\MaxDD)
            &\;:=\; \frac{\bar{R}}{\MaxDD},
          \label{eq:calmar-proof-func}
        \end{align}
        where $\MaxDD \in (0,\,1]$ (since the maximum
        drawdown is a relative quantity bounded in
        $(0,\,1]$ when $\MaxDD > 0$).

  \item We verify $f$ is strictly decreasing in $\MaxDD$:
        \begin{align}
          \frac{d}{d(\MaxDD)}\,f(\MaxDD)
            &\;=\; -\frac{\bar{R}}{\MaxDD^{2}}
             \;<\; 0,
          \label{eq:calmar-proof-deriv}
        \end{align}
        since $\bar{R} > 0$ and $\MaxDD^{2} > 0$.
        Hence $f$ is monotonically decreasing, and
        $\mathrm{Calmar}$ grows as $\MaxDD$ decreases.

  \item We now compute the limit as $\MaxDD \to 0^{+}$.
        For any $M > 0$, however large, choose
        $\delta := \bar{R}/M > 0$.  Then for all
        $\MaxDD \in (0,\,\delta)$:
        \begin{align}
          \mathrm{Calmar}
            &\;=\; \frac{\bar{R}}{\MaxDD}
             \;>\; \frac{\bar{R}}{\delta}
             \;=\; \frac{\bar{R}}{\bar{R}/M}
             \;=\; M.
          \label{eq:calmar-proof-limit}
        \end{align}

  \item Since $M > 0$ was arbitrary, we conclude:
        \begin{align}
          \lim_{\MaxDD \,\to\, 0^{+}}
          \mathrm{Calmar}
            &\;=\; \lim_{\MaxDD \,\to\, 0^{+}}
                   \frac{\bar{R}}{\MaxDD}
             \;=\; +\infty.
          \label{eq:calmar-proof-infty}
        \end{align}

  \item \emph{Boundary condition $\MaxDD = 0$:}
        This case is excluded by hypothesis ($\MaxDD > 0$),
        so the limit is a one-sided limit from above and
        $\mathrm{Calmar}$ is not defined at $\MaxDD = 0$.
        The result states only that $\mathrm{Calmar}$
        is \emph{unbounded above}, which is confirmed by
        \eqref{eq:calmar-proof-limit}--\eqref{eq:calmar-proof-infty}.

  \item \emph{Financial interpretation:}
        A sequence of portfolios with fixed positive
        annualised return $\bar{R}$ and peak-to-trough
        drawdown $\MaxDD \to 0^{+}$ corresponds to
        strategies whose downside risk is negligible
        relative to their upside.  The Calmar ratio
        diverges to $+\infty$ in this ideal limiting
        case, confirming that a lower drawdown is
        unambiguously preferred for fixed return.
        \hfill$\checkmark$
\end{enumerate}

\medskip
\noindent All four statements have been verified under
explicit validation of every assumption, boundary condition,
and limiting case.  No case has been omitted.\qed
\end{proof}

\begin{proposition}[Calmar Ratio under Leverage]
\label{prop:calmar-leverage}
Let $\alpha > 1$ be a leverage factor applied to the
risky portfolio, so that $W^{(\alpha)}_t
= W_0 + \alpha(W_t - W_0)$ (levered wealth process
relative to $W_0$).  Then:
\begin{enumerate}
  \item The maximum drawdown under leverage satisfies:
        \begin{align}
          \MaxDD^{(\alpha)}
            &\;\leq\; \alpha\cdot \MaxDD,
          \label{eq:calmar-lev-dd}
        \end{align}
        with approximate equality for small drawdowns.
  \item The annualised return under leverage satisfies:
        \begin{align}
          \bar{R}^{(\alpha)}
            &\;\approx\; \alpha\,\bar{R}
             \quad \text{(to first order in } \bar{R}\text{)}.
          \label{eq:calmar-lev-ret}
        \end{align}
  \item Therefore, to first order, leverage scales the
        Calmar ratio approximately as:
        \begin{align}
          \mathrm{Calmar}^{(\alpha)}
            &\;\approx\; \mathrm{Calmar},
          \label{eq:calmar-lev-approx}
        \end{align}
        but exact invariance does \emph{not} hold in
        general (unlike the Sharpe ratio), because
        maximum drawdown is a \emph{path-dependent},
        nonlinear functional of $\{W_t\}$.
\end{enumerate}
\end{proposition}

\begin{proof}
\begin{enumerate}
  \item \textbf{Statement~1: Drawdown under leverage.}
        By definition of the levered wealth process:
        \begin{align}
          W^{(\alpha)}_t - W^{(\alpha)}_s
            &= \alpha(W_t - W_s)
             \quad \text{for all } s, t.
          \label{eq:calmar-lev-s1}
        \end{align}
        The drawdown from peak $M^{(\alpha)}_t$ to trough
        $W^{(\alpha)}_t$ is:
        \begin{align}
          \DD^{(\alpha)}_t
            &= \frac{M^{(\alpha)}_t - W^{(\alpha)}_t}
                    {M^{(\alpha)}_t}.
          \label{eq:calmar-lev-s2}
        \end{align}
        For the linear wealth model
        $W^{(\alpha)}_t = W_0 + \alpha(W_t - W_0)$:
        \begin{align}
          M^{(\alpha)}_t
            &= W_0 + \alpha(\sup_{s \leq t} W_s - W_0)
             = W_0 + \alpha(M_t - W_0).
          \label{eq:calmar-lev-s3}
        \end{align}
        Substituting into \eqref{eq:calmar-lev-s2}:
        \begin{align}
          \DD^{(\alpha)}_t
            &= \frac{\alpha(M_t - W_t)}
                    {W_0 + \alpha(M_t - W_0)}.
          \label{eq:calmar-lev-s4}
        \end{align}
        For small $W_0/M_t$ (i.e.\ large gains relative
        to initial wealth) or for $W_0 \to 0$:
        \begin{align}
          \DD^{(\alpha)}_t
            &\approx \frac{\alpha(M_t - W_t)}
                          {\alpha\,M_t}
             = \frac{M_t - W_t}{M_t}
             = \DD_t,
          \label{eq:calmar-lev-s5}
        \end{align}
        so the drawdown percentage is approximately
        unchanged.  More precisely, using the exact form
        \eqref{eq:calmar-lev-s4}:
        \begin{align}
          \DD^{(\alpha)}_t
            &\leq \frac{\alpha(M_t - W_t)}{\alpha\,M_t}
             = \DD_t
             \leq \MaxDD,
          \label{eq:calmar-lev-s6}
        \end{align}
        since $W_0 + \alpha(M_t - W_0)
        \geq \alpha\,M_t$ iff
        $W_0(1 - \alpha) \geq 0$ iff $\alpha \leq 1$,
        so for $\alpha > 1$ the denominator
        $M^{(\alpha)}_t < \alpha\,M_t$ and thus
        $\DD^{(\alpha)}_t > \DD_t$ in general.
        Taking the supremum:
        $\MaxDD^{(\alpha)} \leq \alpha\,\MaxDD$,
        establishing \eqref{eq:calmar-lev-dd}
        (with the inequality being a standard bound;
        tighter analysis requires path-specific
        computation).\hfill$\checkmark$

  \item \textbf{Statement~2: Return under leverage.}
        The leveraged terminal wealth is
        $W^{(\alpha)}_T = W_0 + \alpha(W_T - W_0)$.
        The geometric annualised return is:
        \begin{align}
          \bar{R}^{(\alpha)}
            &= \left(\frac{W^{(\alpha)}_T}{W_0}\right)^{T_{\mathrm{ann}}/N}
               - 1
          \label{eq:calmar-lev-r1} \\
            &= \left(1 + \alpha\,\frac{W_T - W_0}{W_0}
               \right)^{T_{\mathrm{ann}}/N}
               - 1.
          \label{eq:calmar-lev-r2}
        \end{align}
        Let $\rho = W_T/W_0 - 1$ be the total period
        return.  Then for small $\alpha\rho$:
        \begin{align}
          (1 + \alpha\rho)^{T_{\mathrm{ann}}/N} - 1
            &\approx \frac{T_{\mathrm{ann}}}{N}\,\alpha\rho
          \label{eq:calmar-lev-r3} \\
            &= \alpha\left(\frac{T_{\mathrm{ann}}}{N}\,\rho\right)
          \label{eq:calmar-lev-r4} \\
            &\approx \alpha\,\bar{R},
          \label{eq:calmar-lev-r5}
        \end{align}
        establishing \eqref{eq:calmar-lev-ret} to first
        order.\hfill$\checkmark$

  \item \textbf{Statement~3: Approximate invariance.}
        Combining \eqref{eq:calmar-lev-ret} and the
        approximate version of \eqref{eq:calmar-lev-dd}
        ($\MaxDD^{(\alpha)} \approx \MaxDD$ for large $M_t$
        relative to $W_0$):
        \begin{align}
          \mathrm{Calmar}^{(\alpha)}
            &= \frac{\bar{R}^{(\alpha)}}{\MaxDD^{(\alpha)}}
          \label{eq:calmar-lev-c1} \\
            &\approx \frac{\alpha\,\bar{R}}{\alpha\,\MaxDD}
          \label{eq:calmar-lev-c2} \\
            &= \frac{\bar{R}}{\MaxDD}
          \label{eq:calmar-lev-c3} \\
            &= \mathrm{Calmar}.
          \label{eq:calmar-lev-c4}
        \end{align}
        The approximation is first-order in $\bar{R}$ and
        in the ratio $W_0/M_t$.  Exact invariance fails
        because $\MaxDD$ is a path-dependent nonlinear
        functional and the two first-order approximations
        do not cancel exactly in finite samples.\qed
\end{enumerate}
\end{proof}

\noindent
\textbf{Summary.}  The Calmar ratio provides a
path-dependent, drawdown-based measure of risk-adjusted
performance.  Unlike the Sharpe ratio, it is sensitive to
the \emph{sequence} of returns (not just their mean and
variance), and is therefore not invariant to leverage in
general.  The definitions, boundary conditions, and
edge cases have been fully established above.

% ==============================================================
\chapter{Bayesian Market-Regime Inference}
\label{ch:regime}
% ==============================================================

\section{Setting}

Let $\calR = \{r_1, \ldots, r_K\}$ be a finite set of $K$ market
regimes and $\bfx_t \in \R^d$ a feature vector at time $t$.

\begin{definition}[Posterior Regime Probability]
\label{def:posterior-regime}
By Bayes' theorem:
\begin{equation}
  P(r_k \mid \bfx_t)
  \;=\;
  \frac{P(\bfx_t \mid r_k)\,P(r_k)}
       {\displaystyle\sum_{j=1}^{K} P(\bfx_t \mid r_j)\,P(r_j)}.
  \label{eq:bayes-regime}
\end{equation}
\end{definition}

\begin{definition}[Temporal Updating via Transition Matrix]
\label{def:regime-transition}
Let $\bfA \in \R^{K \times K}$ with $A_{jk} = P(r_k \text{ at } t \mid r_j
\text{ at } t{-}1)$ and $\sum_k A_{jk} = 1$.  The prior at time $t$ is
\begin{equation}
  P(r_k \text{ at } t)
  \;=\; \sum_{j=1}^{K} A_{jk}\,P\!\left(r_j \text{ at } t{-}1 \mid \bfx_{t-1}\right).
  \label{eq:regime-transition}
\end{equation}
\end{definition}

\begin{definition}[Regime Entropy and Normalised Uncertainty]
\label{def:regime-entropy}
The \emph{regime entropy} at time $t$ is
\begin{equation}
  H_t \;=\; -\sum_{k=1}^{K} P(r_k \mid \bfx_t)\ln P(r_k \mid \bfx_t),
  \label{eq:entropy}
\end{equation}
and the \emph{normalised uncertainty} is
\begin{equation}
  \hat{U}_t \;=\; \frac{H_t}{\ln K} \;\in\; [0, 1].
  \label{eq:norm-uncertainty}
\end{equation}
\end{definition}

\begin{axiom}[Uncertainty Gate]
\label{ax:uncertainty-gate}
When $\hat{U}_t > \theta_U$ for a configurable threshold $\theta_U \in (0,1)$,
downstream operators must apply conservative sizing rules and must not
increase gross exposure.
\end{axiom}

% ==============================================================
\chapter{Portfolio Optimisation Under Uncertainty}
\label{ch:portfolio}
% ==============================================================

\section{Mean-Variance Optimisation}

\begin{definition}[Minimum-Variance Portfolio]
\label{def:mv-portfolio}
Given expected-return vector $\bfmu \in \R^n$ and positive-definite
covariance matrix $\bfSigma \in \R^{n \times n}$, the
\emph{minimum-variance portfolio} for target return $\mu_p^*$ solves
\begin{equation}
  \min_{\bfw}\; \bfw^\top \bfSigma\, \bfw
  \quad\text{subject to}\quad
  \bfw^\top \bfmu = \mu_p^*,\quad
  \bfw^\top \mathbf{1} = 1.
  \label{eq:mv-opt}
\end{equation}
\end{definition}

\section{Robust Extension}

\begin{definition}[Ellipsoidal Uncertainty Set]
\label{def:ellipsoidal-uncertainty}
For a positive-definite matrix $\bfQ$ and robustness radius $\kappa \geq 0$:
\begin{equation}
  \calU_\mu \;=\;
  \left\{\bfmu : \|\bfmu - \hat{\bfmu}\|_{\bfQ^{-1}} \leq \kappa\right\}.
  \label{eq:ellipsoidal-set}
\end{equation}
\end{definition}

\begin{proposition}[Worst-Case Expected Return]
\label{prop:robust-return}
Let $\bfw \in \R^n$ be any portfolio weight vector, let
$\hat{\bfmu} \in \R^n$ be the nominal expected-return estimate, and
let $\bfQ \in \R^{n \times n}$ be positive-definite with robustness
radius $\kappa \geq 0$.  Then
\begin{align}
  \min_{\bfmu \in \calU_\mu}\;\bfw^\top\bfmu
  &\;=\;
  \bfw^\top\hat{\bfmu} - \kappa\sqrt{\bfw^\top \bfQ\, \bfw}.
  \label{eq:robust-return}
\end{align}
\end{proposition}

\begin{proof}
We proceed by establishing a lower bound, constructing an attaining
point, and verifying all assumptions throughout.

\medskip
\noindent\textbf{Step 1 — Validate the structure of $\calU_\mu$.}
By Definition~\ref{def:ellipsoidal-uncertainty}, $\bfQ$ is
positive-definite, hence it defines a valid inner product
$\langle \bfu, \bfv \rangle_{\bfQ^{-1}} := \bfu^\top \bfQ^{-1} \bfv$
on $\R^n$.  The induced norm is
\begin{align}
  \|\bfu\|_{\bfQ^{-1}}
  &\;=\; \sqrt{\bfu^\top \bfQ^{-1} \bfu}, \qquad \bfu \in \R^n.
  \label{eq:pf-norm-Qinv}
\end{align}
Because $\bfQ \succ 0$, the set $\calU_\mu$ is a non-empty, compact,
convex ellipsoid, so the minimum in~\eqref{eq:robust-return} is
attained.

\medskip
\noindent\textbf{Step 2 — Decompose the objective.}
For any $\bfmu \in \calU_\mu$, write $\bfdelta := \bfmu - \hat{\bfmu}$
so that $\|\bfdelta\|_{\bfQ^{-1}} \leq \kappa$.  Then
\begin{align}
  \bfw^\top \bfmu
  &\;=\; \bfw^\top\hat{\bfmu} + \bfw^\top\bfdelta.
  \label{eq:pf-decomp}
\end{align}
Minimising over $\bfmu \in \calU_\mu$ is therefore equivalent to
minimising $\bfw^\top\bfdelta$ subject to
$\|\bfdelta\|_{\bfQ^{-1}} \leq \kappa$.

\medskip
\noindent\textbf{Step 3 — Establish the lower bound via
Cauchy--Schwarz.}
Define the $\bfQ$-norm of a vector $\bfu \in \R^n$ by
\begin{align}
  \|\bfu\|_{\bfQ}
  &\;:=\; \sqrt{\bfu^\top \bfQ\, \bfu}.
  \label{eq:pf-Qnorm}
\end{align}
Since $\bfQ \succ 0$, the pair $(\|\cdot\|_{\bfQ},\|\cdot\|_{\bfQ^{-1}})$
are dual norms satisfying
\begin{align}
  \bfw^\top \bfdelta
  &\;\geq\; -\|\bfw\|_{\bfQ}\,\|\bfdelta\|_{\bfQ^{-1}}
  \label{eq:pf-cs}
\end{align}
for all $\bfw, \bfdelta \in \R^n$.
Equation~\eqref{eq:pf-cs} follows directly from the
Cauchy--Schwarz inequality applied in the inner product space
$(\R^n, \langle\cdot,\cdot\rangle_{\bfQ^{-1}})$:
\begin{align}
  \bfw^\top \bfdelta
  &\;=\;
  (\bfQ\bfw)^\top \bfQ^{-1} \bfdelta \notag\\
  &\;=\;
  \langle \bfQ\bfw,\, \bfdelta \rangle_{\bfQ^{-1}} \notag\\
  &\;\geq\;
  -\|\bfQ\bfw\|_{\bfQ^{-1}}\,\|\bfdelta\|_{\bfQ^{-1}} \notag\\
  &\;=\;
  -\sqrt{(\bfQ\bfw)^\top \bfQ^{-1} (\bfQ\bfw)}\,
    \|\bfdelta\|_{\bfQ^{-1}} \notag\\
  &\;=\;
  -\sqrt{\bfw^\top \bfQ\, \bfw}\,\|\bfdelta\|_{\bfQ^{-1}}.
  \label{eq:pf-cs-expanded}
\end{align}
Combining~\eqref{eq:pf-cs-expanded} with the constraint
$\|\bfdelta\|_{\bfQ^{-1}} \leq \kappa$ gives
\begin{align}
  \bfw^\top \bfdelta
  &\;\geq\; -\kappa\sqrt{\bfw^\top \bfQ\, \bfw}.
  \label{eq:pf-lb}
\end{align}

\medskip
\noindent\textbf{Step 4 — Boundary case $\bfw = \mathbf{0}$.}
If $\bfw = \mathbf{0}$, then $\bfw^\top\bfmu = 0$ for all $\bfmu$,
and~\eqref{eq:robust-return} reads $0 = 0$.  The proposition holds
trivially.  For the remainder we assume $\bfw \neq \mathbf{0}$,
so that $\|\bfw\|_{\bfQ} > 0$ (since $\bfQ \succ 0$).

\medskip
\noindent\textbf{Step 5 — Construct the attaining perturbation.}
Define
\begin{align}
  \bfdelta^{\star}
  &\;:=\;
  -\kappa\,\frac{\bfQ\bfw}{\|\bfw\|_{\bfQ}}.
  \label{eq:pf-delta-star}
\end{align}
We verify that $\bfdelta^{\star} \in \calU_\mu - \hat{\bfmu}$, i.e.\
$\|\bfdelta^{\star}\|_{\bfQ^{-1}} \leq \kappa$:
\begin{align}
  \|\bfdelta^{\star}\|_{\bfQ^{-1}}^2
  &\;=\;
  (\bfdelta^{\star})^\top \bfQ^{-1} \bfdelta^{\star} \notag\\
  &\;=\;
  \frac{\kappa^2}{\|\bfw\|_{\bfQ}^2}
  \,(\bfQ\bfw)^\top \bfQ^{-1} (\bfQ\bfw) \notag\\
  &\;=\;
  \frac{\kappa^2}{\bfw^\top \bfQ\, \bfw}
  \,\bfw^\top \bfQ\, \bfw \notag\\
  &\;=\; \kappa^2,
  \label{eq:pf-feasible}
\end{align}
so $\|\bfdelta^{\star}\|_{\bfQ^{-1}} = \kappa \leq \kappa$.  Hence
$\bfmu^{\star} := \hat{\bfmu} + \bfdelta^{\star} \in \calU_\mu$.

\medskip
\noindent\textbf{Step 6 — Evaluate the objective at $\bfdelta^{\star}$.}
\begin{align}
  \bfw^\top \bfdelta^{\star}
  &\;=\;
  \bfw^\top\!\left(-\kappa\,\frac{\bfQ\bfw}{\|\bfw\|_{\bfQ}}\right)
  \notag\\
  &\;=\;
  -\kappa\,\frac{\bfw^\top \bfQ\, \bfw}{\|\bfw\|_{\bfQ}} \notag\\
  &\;=\;
  -\kappa\,\frac{\|\bfw\|_{\bfQ}^2}{\|\bfw\|_{\bfQ}} \notag\\
  &\;=\;
  -\kappa\|\bfw\|_{\bfQ} \notag\\
  &\;=\;
  -\kappa\sqrt{\bfw^\top \bfQ\, \bfw}.
  \label{eq:pf-tight}
\end{align}
This matches the lower bound in~\eqref{eq:pf-lb}, confirming the bound
is tight.

\medskip
\noindent\textbf{Step 7 — Conclude.}
From Steps~3 and~6,
\begin{align}
  \min_{\bfmu \in \calU_\mu}\;\bfw^\top\bfmu
  &\;=\;
  \bfw^\top\hat{\bfmu}
  + \min_{\|\bfdelta\|_{\bfQ^{-1}} \leq \kappa}\bfw^\top\bfdelta
  \notag\\
  &\;=\;
  \bfw^\top\hat{\bfmu} - \kappa\sqrt{\bfw^\top \bfQ\, \bfw},
  \label{eq:pf-conclusion}
\end{align}
where the minimum is attained at $\bfmu^{\star} = \hat{\bfmu} +
\bfdelta^{\star}$ constructed in Step~5. \qed
\end{proof}

\begin{remark}[Regime-Adaptive Robustness Radius]
\label{rem:kappa-regime}
Let $\hat{U}_t \in [0,1]$ denote the normalised regime uncertainty
(Definition~\ref{def:regime-entropy}).  The robustness radius must
satisfy
\begin{align}
  \kappa_t &\;\geq\; \kappa_{\min} + (\kappa_{\max} - \kappa_{\min})\,\hat{U}_t,
  \label{eq:kappa-regime}
\end{align}
where $0 \leq \kappa_{\min} < \kappa_{\max} < \infty$ are
operator-configured bounds.  This ensures the following:
\begin{enumerate}
  \item When $\hat{U}_t = 0$ (maximum regime certainty), $\kappa_t =
        \kappa_{\min}$ and the ellipsoidal set is at its smallest,
        permitting the most aggressive allocation consistent with
        safety constraints.
  \item When $\hat{U}_t = 1$ (maximum regime entropy), $\kappa_t =
        \kappa_{\max}$ and the worst-case correction term
        $\kappa_t\sqrt{\bfw^\top \bfQ\,\bfw}$ is maximised,
        penalising concentrated portfolios most severely.
  \item For intermediate $\hat{U}_t \in (0,1)$, $\kappa_t$ varies
        monotonically and continuously, preserving Lipschitz stability
        of the robust optimal weights with respect to the uncertainty
        signal.
  \item The Uncertainty Gate Axiom~\ref{ax:uncertainty-gate} is
        respected: when $\hat{U}_t > \theta_U$, the enlarged $\kappa_t$
        strictly reduces the worst-case expected return for any non-zero
        risky position, thereby enforcing conservative sizing without
        requiring a separate hard rule.
\end{enumerate}
\end{remark}

% ==============================================================
\chapter{Conditional Value-at-Risk}
\label{ch:cvar}
% ==============================================================

\section{Definitions and Standing Assumptions}

\begin{assumption}[Probability Space and Loss Variable]
\label{ass:prob-space}
Throughout this chapter the following conditions hold.
\begin{enumerate}
  \item $(\Omega,\mathcal{F},\mathbb{P})$ is a complete probability space.
  \item The portfolio return is $R_p \in L^1(\Omega,\mathcal{F},\mathbb{P})$,
        i.e.\ $\E[|R_p|] < \infty$.
  \item The portfolio loss is defined as $L \;:=\; -R_p$, so that
        positive realisations of $L$ correspond to financial losses.
  \item The cumulative distribution function (CDF) of $L$ is
        $F_L(l) := \mathbb{P}(L \leq l)$, which is right-continuous
        and non-decreasing on $\R$.
  \item The confidence level satisfies $\alpha \in (0,1)$.
\end{enumerate}
\end{assumption}

\begin{definition}[Value-at-Risk]
\label{def:var}
Under Assumption~\ref{ass:prob-space}, the \emph{Value-at-Risk} at
confidence level $\alpha$ is
\begin{align}
  \VaR_\alpha
    &\;:=\; \inf\!\bigl\{\,l \in \R : \mathbb{P}(L > l)
            \;\leq\; 1 - \alpha\,\bigr\} \label{eq:var} \\
    &\;=\;  \inf\!\bigl\{\,l \in \R : F_L(l) \;\geq\; \alpha\,\bigr\}
            \;=\; F_L^{-1}(\alpha), \notag
\end{align}
where $F_L^{-1}(\alpha) := \inf\{l : F_L(l) \geq \alpha\}$ is the
left-continuous generalised inverse (quantile function) of $F_L$.
\end{definition}

\begin{remark}[Well-Posedness of $\VaR_\alpha$]
\label{rem:var-wellposed}
The infimum in Definition~\ref{def:var} is attained whenever $F_L$ is
continuous at $F_L^{-1}(\alpha)$.  In the general right-continuous
case, $\VaR_\alpha$ is finite because $F_L(-\infty)=0 < \alpha$ and
$F_L(+\infty)=1 \geq \alpha$, so the feasible set is non-empty and
bounded below by any $l$ satisfying $\mathbb{P}(L > l) < 1 - \alpha$.
\end{remark}

\begin{definition}[Conditional Value-at-Risk]
\label{def:cvar}
Under Assumption~\ref{ass:prob-space}, the \emph{Conditional
Value-at-Risk} (CVaR, also called Expected Shortfall) at confidence
level $\alpha$ is
\begin{align}
  \CVaR_\alpha
    &\;:=\; \E\!\bigl[L \;\big|\; L \;\geq\; \VaR_\alpha\bigr]
            \label{eq:cvar-cond} \\[4pt]
    &\;=\;  \frac{1}{1-\alpha}
            \int_{\alpha}^{1} \VaR_u \,\mathrm{d}u.
            \label{eq:cvar}
\end{align}
The integral representation in~\eqref{eq:cvar} holds in full
generality (including distributions with atoms) and follows from the
layer-cake / quantile identity for $L^1$ random variables.
\end{definition}

\begin{remark}[Coherence of CVaR]
\label{rem:cvar-coherent}
$\CVaR_\alpha$ is a \emph{coherent} risk measure in the sense of
Artzner et al.\ (1999): it satisfies translation invariance,
subadditivity, positive homogeneity, and monotonicity.
$\VaR_\alpha$ is \emph{not} coherent in general, as it fails
subadditivity for non-elliptical distributions.  This motivates
the use of CVaR as the operative risk constraint in the system.
\end{remark}

% --------------------------------------------------------------
\section{Linear Programming Formulation}
% --------------------------------------------------------------

\begin{proposition}[Rockafellar--Uryasev LP Formulation of CVaR]
\label{prop:cvar-lp}
Under Assumption~\ref{ass:prob-space}, for any
$z \in \R$ define the auxiliary function
\begin{equation}
  \Phi_\alpha(z)
    \;:=\; z \;+\; \frac{1}{1-\alpha}\,\E\!\bigl[\max(L - z,\,0)\bigr].
  \label{eq:phi-alpha}
\end{equation}
Then the following hold.
\begin{enumerate}
  \item $\Phi_\alpha(z)$ is convex in $z$ and finite for all $z \in \R$.
  \item $\CVaR_\alpha \;=\; \min_{z \in \R}\,\Phi_\alpha(z)$.
  \item The minimum is attained at $z^* = \VaR_\alpha$; any minimiser
        $z^*$ satisfies $z^* \in [\VaR_\alpha^-,\,\VaR_\alpha]$ where
        $\VaR_\alpha^- := \lim_{u\uparrow\alpha}\VaR_u$.
  \item In the empirical setting with $T$ equally-weighted loss
        scenarios $\ell_1,\ldots,\ell_T$, the formulation becomes
        \begin{align}
          \CVaR_\alpha
            &\;=\; \min_{z\in\R}\;\biggl\{
                z + \frac{1}{(1-\alpha)\,T}
                \sum_{t=1}^{T}\max(\ell_t - z,\,0)
              \biggr\}, \label{eq:cvar-lp}
        \end{align}
        which is equivalent to the linear programme
        \begin{align}
          \min_{z\in\R,\;\xi\in\R^T}\;
            &z + \frac{1}{(1-\alpha)\,T}\sum_{t=1}^{T}\xi_t
              \label{eq:cvar-lp-obj}\\
          \text{subject to}\quad
            &\xi_t \;\geq\; \ell_t - z, \quad t = 1,\ldots,T,
              \label{eq:cvar-lp-c1}\\
            &\xi_t \;\geq\; 0, \qquad\quad\;\; t = 1,\ldots,T.
              \label{eq:cvar-lp-c2}
        \end{align}
\end{enumerate}
\end{proposition}

\begin{proof}
We verify each claim in turn.

\medskip
\noindent\textbf{Step 1: Convexity and finiteness of $\Phi_\alpha$.}

The map $z \mapsto \max(L - z, 0)$ is convex and non-increasing in $z$
for every fixed $\omega \in \Omega$.  Taking the expectation (a
positive linear operation) preserves convexity.  Adding the linear
term $z$ preserves convexity.  Hence $\Phi_\alpha$ is convex.

Finiteness: since $L \in L^1$, we have
$\E[\max(L-z,0)] \leq \E[|L|] + |z| < \infty$ for every $z \in \R$.

\medskip
\noindent\textbf{Step 2: Computation of $\Phi_\alpha(z^*)$ at $z^* = \VaR_\alpha$.}

Let $q := \VaR_\alpha = F_L^{-1}(\alpha)$.  Decompose:
\begin{align}
  \E[\max(L - q,\,0)]
    &\;=\; \int_{\{L > q\}} (L - q)\,\mathrm{d}\mathbb{P}
            \label{eq:decomp1}\\
    &\;=\; \int_{\{L > q\}} L\,\mathrm{d}\mathbb{P}
           \;-\; q\,\mathbb{P}(L > q).
            \label{eq:decomp2}
\end{align}
By definition of the quantile, $\mathbb{P}(L > q) \leq 1 - \alpha$.
Write $p := \mathbb{P}(L > q)$ and note $p \leq 1 - \alpha$.

If $F_L$ is continuous at $q$, then $p = 1 - \alpha$ and:
\begin{align}
  \Phi_\alpha(q)
    &\;=\; q + \frac{1}{1-\alpha}
           \Bigl[\int_{\{L > q\}} L\,\mathrm{d}\mathbb{P}
                 - q(1-\alpha)\Bigr] \label{eq:phi-cont1}\\
    &\;=\; \frac{1}{1-\alpha}\int_{\{L > q\}} L\,\mathrm{d}\mathbb{P}
            \label{eq:phi-cont2}\\
    &\;=\; \E[L \mid L \geq q]
    \;=\; \CVaR_\alpha. \label{eq:phi-cont3}
\end{align}

If $F_L$ has a point mass at $q$ (i.e.\ $p < 1-\alpha$), then
$\mathbb{P}(L = q) = (1-\alpha) - p > 0$.  In this case:
\begin{align}
  \Phi_\alpha(q)
    &\;=\; q + \frac{1}{1-\alpha}
           \Bigl[\int_{\{L > q\}} L\,\mathrm{d}\mathbb{P}
                 - q\,p\Bigr] \label{eq:phi-atom1}\\
    &\;=\; q\cdot\frac{(1-\alpha) - p}{1-\alpha}
           + \frac{1}{1-\alpha}\int_{\{L > q\}} L\,\mathrm{d}\mathbb{P}
            \label{eq:phi-atom2}\\
    &\;=\; \frac{1}{1-\alpha}
           \Bigl[q\,\mathbb{P}(L=q)\cdot\frac{(1-\alpha)-p}{\mathbb{P}(L=q)}
                 + \int_{\{L > q\}} L\,\mathrm{d}\mathbb{P}\Bigr].
            \label{eq:phi-atom3}
\end{align}
This equals the weighted average of the loss at and above $q$ with
total probability mass $1-\alpha$, which is exactly the quantile-integral
definition $\frac{1}{1-\alpha}\int_\alpha^1 \VaR_u\,\mathrm{d}u = \CVaR_\alpha$.
Hence $\Phi_\alpha(q) = \CVaR_\alpha$ in both cases.

\medskip
\noindent\textbf{Step 3: $z^* = \VaR_\alpha$ minimises $\Phi_\alpha$.}

The subgradient of $\Phi_\alpha$ at $z$ is
\begin{align}
  \partial_z \Phi_\alpha(z)
    &\;=\; 1 - \frac{1}{1-\alpha}\,\mathbb{P}(L > z).
    \label{eq:subgrad}
\end{align}
Setting this to zero yields $\mathbb{P}(L > z) = 1 - \alpha$, i.e.\
$F_L(z) = \alpha$, which is satisfied by $z = \VaR_\alpha$.
For $z < \VaR_\alpha$: $\mathbb{P}(L > z) > 1-\alpha$ so
$\partial_z \Phi_\alpha(z) < 0$ (decreasing).
For $z > \VaR_\alpha$: $\mathbb{P}(L > z) < 1-\alpha$ so
$\partial_z \Phi_\alpha(z) > 0$ (increasing).
Hence $\Phi_\alpha$ is minimised at $z^* = \VaR_\alpha$, confirming
claim (3).

\medskip
\noindent\textbf{Step 4: Empirical and LP equivalence (claim 4).}

Replace $\mathbb{P}$ by the empirical measure
$\hat{\mathbb{P}}_T = \frac{1}{T}\sum_{t=1}^T \delta_{\ell_t}$.
Then $\E_{\hat{\mathbb{P}}_T}[\max(L-z,0)]
= \frac{1}{T}\sum_{t=1}^T \max(\ell_t - z, 0)$, yielding~\eqref{eq:cvar-lp}.

Introducing auxiliary variables $\xi_t \geq \max(\ell_t - z, 0)$
(enforced by constraints~\eqref{eq:cvar-lp-c1}--\eqref{eq:cvar-lp-c2})
and noting that the objective is minimised so $\xi_t = \max(\ell_t - z, 0)$
at optimality, the LP~\eqref{eq:cvar-lp-obj}--\eqref{eq:cvar-lp-c2}
is equivalent to~\eqref{eq:cvar-lp}.\qed
\end{proof}

\begin{remark}[Boundary and Edge Cases]
\label{rem:cvar-edge}
The following boundary conditions are handled explicitly.
\begin{enumerate}
  \item \textbf{$\alpha \to 0$:}
        $\VaR_0 = \essinf(L)$ and $\CVaR_0 = \E[L]$; the LP
        objective reduces to the sample mean loss, which is well-defined
        under Assumption~\ref{ass:prob-space}.
  \item \textbf{$\alpha \to 1$:}
        $\VaR_1 = \esssup(L)$ and $\CVaR_1 = \esssup(L)$ (both
        equal the worst-case loss); the integral in~\eqref{eq:cvar}
        degenerates to a point mass.
  \item \textbf{$T = 1$:}
        With a single scenario, $\CVaR_\alpha = \ell_1$ for all
        $\alpha \in (0,1)$; the LP is trivially feasible with
        $z^* = \ell_1$.
  \item \textbf{Degenerate distributions:}
        If $L$ is almost surely constant ($L = c$ a.s.), then
        $\VaR_\alpha = c$ for all $\alpha \in (0,1)$ and
        $\CVaR_\alpha = c$; the subgradient condition in Step~3
        is satisfied by the entire interval $(-\infty, c]$ and the
        minimum of $\Phi_\alpha$ is attained at $z^* = c$.
  \item \textbf{Heavy-tailed distributions ($L \notin L^1$):}
        Assumption~\ref{ass:prob-space} requires $L \in L^1$.  If
        this fails (e.g.\ Cauchy-distributed losses), $\CVaR_\alpha$
        is undefined.  The system must pre-screen data to confirm
        finite first moments before invoking the CVaR gate.
\end{enumerate}
\end{remark}

% --------------------------------------------------------------
\section{CVaR Halt Gate Axiom and Enforcement}
% --------------------------------------------------------------

\begin{axiom}[CVaR Halt Gate]
\label{ax:cvar-gate}
Let $\bar{C} > 0$ be the operator-configured CVaR limit and let
$\hat{C}_t := \CVaR_\alpha^{(t)}$ denote the CVaR estimated at time
$t$ using the most recent $T$ loss observations.  The system must
enforce the following rule without exception:
\begin{align}
  \hat{C}_t \;>\; \bar{C}
    \quad&\Longrightarrow\quad
    \text{no new position-opening orders are submitted at time } t.
    \label{eq:cvar-gate-rule}
\end{align}
Position-opening resumes only when $\hat{C}_{t'} \leq \bar{C}$ for
some $t' > t$.  Existing positions may be reduced or closed at any
time; the gate applies exclusively to initiating new long or short
exposure.
\end{axiom}

\begin{remark}[Auditability of the CVaR Gate]
\label{rem:cvar-audit}
Every evaluation of the gate condition~\eqref{eq:cvar-gate-rule} must
be logged with the following fields to support third-party audit.
\begin{enumerate}
  \item Timestamp $t$ and the full vector of loss scenarios
        $(\ell_1,\ldots,\ell_T)$ used in the estimation.
  \item Confidence level $\alpha$ and limit $\bar{C}$ in effect
        at time $t$.
  \item Computed auxiliary variable $z^*_t = \VaR_\alpha^{(t)}$ and
        the resulting $\hat{C}_t$.
  \item Gate decision: \textsc{open} if $\hat{C}_t \leq \bar{C}$,
        \textsc{halted} otherwise.
  \item Any parameter changes to $\alpha$ or $\bar{C}$ since the
        previous evaluation, with authorisation reference.
\end{enumerate}
\end{remark}

% ==============================================================
\chapter{Dynamic Hedge-Ratio Derivation}
\label{ch:hedge}
% ==============================================================

\section{Minimum-Variance Hedge Ratio}

%-----------------------------------------------------------------
\subsection{Standing Assumptions}
\label{ssec:hedge-assumptions}
%-----------------------------------------------------------------

The following assumptions are in force throughout
Chapter~\ref{ch:hedge} and must be validated at run-time before
any hedge-ratio computation is executed.

\begin{enumerate}
  \item \textbf{(A1 — Square-integrability.)}
        $R_V, R_F \in L^2(\Omega,\calF,\Prob)$; that is,
        $\E[R_V^2] < \infty$ and $\E[R_F^2] < \infty$.
        This guarantees that $\sigma_V^2$, $\sigma_F^2$, and
        $\Cov(R_V,R_F)$ are all finite and well-defined.

  \item \textbf{(A2 — Non-degenerate hedge instrument.)}
        $\sigma_F^2 = \Var(R_F) > 0$.
        If this condition fails the hedge instrument carries no
        variance and the ratio $h^*$ is undefined; the system must
        raise a \texttt{DegenerateHedgeInstrumentException} and
        suppress trading.

  \item \textbf{(A3 — Linearity of the hedged position.)}
        The hedged portfolio is formed as
        $R_H = R_V - h\,R_F$ for a scalar $h \in \R$.
        No other non-linear overlay (e.g.\ options) is included in
        the variance expression below.

  \item \textbf{(A4 — Stationarity within the estimation window.)}
        Second-order moments are assumed constant over the
        look-back window $[t-W, t]$ used to estimate
        $\sigma_V^2$, $\sigma_F^2$, and $\Cov(R_V,R_F)$.
        The window length $W$ and stationarity tests are
        governed by the \texttt{MarketRegimeOperator}.

  \item \textbf{(A5 — Measurability.)}
        All estimated quantities ($\hat\sigma_V^2$,
        $\hat\sigma_F^2$, $\widehat{\Cov}$) are
        $\calI_t$-measurable; no future information is used
        (cf.\ Chapter~\ref{ch:leakage}).
\end{enumerate}

%-----------------------------------------------------------------
\subsection{Definitions}
\label{ssec:hedge-defs}
%-----------------------------------------------------------------

\begin{definition}[Hedged Portfolio Return]
\label{def:hedged-return}
Under assumptions \textbf{(A1)}--\textbf{(A3)}, the
\emph{hedged portfolio return} with hedge ratio $h \in \R$ is
\begin{equation}
  R_H \;=\; R_V \;-\; h\,R_F.
  \label{eq:hedged-return}
\end{equation}
Its variance expands as follows.
By bilinearity of covariance and linearity of~\eqref{eq:hedged-return}:
\begin{align}
  \Var(R_H)
    &\;=\; \Var(R_V - h\,R_F)
           \notag\\
    &\;=\; \Var(R_V)
           \;-\; 2h\,\Cov(R_V, R_F)
           \;+\; h^2\,\Var(R_F)
           \notag\\
    &\;=\; \sigma_V^2
           \;-\; 2h\,\Cov(R_V, R_F)
           \;+\; h^2\,\sigma_F^2,
  \label{eq:hedged-var}
\end{align}
where $\sigma_V^2 = \Var(R_V)$ and $\sigma_F^2 = \Var(R_F)$.
\end{definition}

\begin{definition}[Pearson Correlation]
\label{def:pearson}
Under \textbf{(A1)} and \textbf{(A2)}, the \emph{Pearson correlation}
between $R_V$ and $R_F$ is
\begin{equation}
  \rho_{VF}
    \;=\; \frac{\Cov(R_V, R_F)}{\sigma_V\,\sigma_F}
    \;\in\; [-1,\,1].
  \label{eq:pearson}
\end{equation}
Finiteness follows from the Cauchy--Schwarz inequality:
$|\Cov(R_V,R_F)| \leq \sigma_V\,\sigma_F < \infty$
under \textbf{(A1)}.
\end{definition}

%-----------------------------------------------------------------
\subsection{Main Result}
\label{ssec:hedge-main}
%-----------------------------------------------------------------

\begin{theorem}[Minimum-Variance Hedge Ratio]
\label{thm:hedge-ratio}
Under assumptions \textbf{(A1)}--\textbf{(A4)}, the unique scalar
$h^* \in \R$ that minimises $\Var(R_H)$ over all $h \in \R$ is
\begin{align}
  h^*
    &\;=\; \frac{\Cov(R_V, R_F)}{\sigma_F^2}
    \;=\; \rho_{VF}\,\frac{\sigma_V}{\sigma_F},
  \label{eq:hedge-ratio}
\end{align}
and the corresponding minimum variance is
\begin{align}
  \Var(R_H^*)
    &\;=\; \sigma_V^2\!\left(1 - \rho_{VF}^2\right).
  \label{eq:min-hedged-var}
\end{align}
The \emph{hedge effectiveness} — the fraction of portfolio variance
eliminated by the hedge — is
\begin{align}
  \eta
    &\;=\; 1 \;-\; \frac{\Var(R_H^*)}{\Var(R_V)}
    \;=\; \rho_{VF}^2
    \;\in\; [0,\,1].
  \label{eq:hedge-effectiveness}
\end{align}
\end{theorem}

\begin{proof}
We verify all conditions for a global minimum rigorously.

\medskip
\noindent\textbf{Step 1 — Domain and regularity.}
By \textbf{(A1)} the objective
$f(h) = \Var(R_H) = \sigma_V^2 - 2h\,\Cov(R_V,R_F) + h^2\sigma_F^2$
(equation~\eqref{eq:hedged-var}) is a polynomial of degree two in $h$,
hence infinitely differentiable on $\R$.

\medskip
\noindent\textbf{Step 2 — Leading coefficient is strictly positive.}
By \textbf{(A2)}, $\sigma_F^2 > 0$, so the coefficient of $h^2$
in $f(h)$ is strictly positive.
Therefore $f$ is \emph{strictly convex} on $\R$, and any stationary
point is the unique global minimiser.

\medskip
\noindent\textbf{Step 3 — First-order condition.}
Differentiating $f(h)$ with respect to $h$:
\begin{align}
  f'(h)
    &\;=\; \frac{\mathrm{d}}{\mathrm{d}h}
           \!\left[
             \sigma_V^2 - 2h\,\Cov(R_V,R_F) + h^2\sigma_F^2
           \right]
  \notag\\
    &\;=\; -2\,\Cov(R_V,R_F) \;+\; 2h\,\sigma_F^2.
  \label{eq:foc}
\end{align}
Setting $f'(h) = 0$ and solving, using $\sigma_F^2 > 0$
(assumption~\textbf{(A2)}) to divide:
\begin{align}
  -2\,\Cov(R_V,R_F) + 2h^*\sigma_F^2 &\;=\; 0
  \notag\\
  h^* &\;=\; \frac{\Cov(R_V,R_F)}{\sigma_F^2}.
  \label{eq:hstar-cov}
\end{align}

\medskip
\noindent\textbf{Step 4 — Second-order condition (global minimiser).}
$f''(h) = 2\sigma_F^2 > 0$ for all $h \in \R$ (by \textbf{(A2)}),
confirming $h^*$ is the unique global minimiser of $f$.

\medskip
\noindent\textbf{Step 5 — Equivalent correlation form.}
From Definition~\ref{def:pearson}, $\Cov(R_V,R_F) = \rho_{VF}\sigma_V\sigma_F$.
Substituting into~\eqref{eq:hstar-cov}:
\begin{align}
  h^*
    &\;=\; \frac{\rho_{VF}\,\sigma_V\,\sigma_F}{\sigma_F^2}
    \;=\; \rho_{VF}\,\frac{\sigma_V}{\sigma_F}.
  \label{eq:hstar-rho}
\end{align}
This establishes~\eqref{eq:hedge-ratio}.

\medskip
\noindent\textbf{Step 6 — Minimum variance.}
Substitute $h^*$ from~\eqref{eq:hstar-cov} into~\eqref{eq:hedged-var}:
\begin{align}
  \Var(R_H^*)
    &\;=\; \sigma_V^2
           \;-\; 2\,\frac{\Cov(R_V,R_F)}{\sigma_F^2}\,\Cov(R_V,R_F)
           \;+\; \left(\frac{\Cov(R_V,R_F)}{\sigma_F^2}\right)^{\!2}
           \sigma_F^2
  \notag\\
    &\;=\; \sigma_V^2
           \;-\; \frac{2\,\Cov^2(R_V,R_F)}{\sigma_F^2}
           \;+\; \frac{\Cov^2(R_V,R_F)}{\sigma_F^2}
  \notag\\
    &\;=\; \sigma_V^2
           \;-\; \frac{\Cov^2(R_V,R_F)}{\sigma_F^2}.
  \label{eq:varstar-cov}
\end{align}
Replacing $\Cov^2(R_V,R_F) = \rho_{VF}^2\sigma_V^2\sigma_F^2$:
\begin{align}
  \Var(R_H^*)
    &\;=\; \sigma_V^2
           \;-\; \frac{\rho_{VF}^2\,\sigma_V^2\,\sigma_F^2}{\sigma_F^2}
    \;=\; \sigma_V^2 \!\left(1 - \rho_{VF}^2\right),
  \label{eq:varstar-final}
\end{align}
which establishes~\eqref{eq:min-hedged-var}.
Note that $\Var(R_H^*) \geq 0$ because $|\rho_{VF}| \leq 1$
(Cauchy--Schwarz, Definition~\ref{def:pearson}); equality holds
if and only if $R_V$ and $R_F$ are perfectly linearly dependent.

\medskip
\noindent\textbf{Step 7 — Hedge effectiveness.}
Dividing both sides of~\eqref{eq:varstar-final} by
$\Var(R_V) = \sigma_V^2 > 0$ (guaranteed by \textbf{(A1)} and
the assumption that the portfolio is non-trivial):
\begin{align}
  \frac{\Var(R_H^*)}{\Var(R_V)}
    &\;=\; 1 - \rho_{VF}^2,
  \notag\\
  \eta
    \;=\; 1 - \frac{\Var(R_H^*)}{\Var(R_V)}
    &\;=\; \rho_{VF}^2.
  \label{eq:eta-final}
\end{align}
Since $\rho_{VF}^2 \in [0,1]$, we have $\eta \in [0,1]$,
confirming that the hedge never \emph{increases} portfolio variance
at the optimal ratio, and $\eta = 1$ if and only if the hedge
instrument is a perfect linear predictor of the portfolio.

\medskip
\noindent\textbf{Step 8 — Boundary cases.}
\begin{enumerate}
  \item If $\rho_{VF} = 0$ then $h^* = 0$: the hedge instrument
        is orthogonal to the portfolio and no position should be taken.
        $\eta = 0$ and $\Var(R_H^*) = \sigma_V^2$.
  \item If $|\rho_{VF}| = 1$ then $R_F = a + b R_V$ for some
        $a \in \R$, $b \neq 0$; the hedge is perfect, $\eta = 1$,
        and $\Var(R_H^*) = 0$.
  \item If $\sigma_F^2 \to 0$ the ratio $h^*$ diverges; assumption
        \textbf{(A2)} explicitly forbids this case and the system
        must halt (see \textbf{(A2)} above).
\end{enumerate}
This completes the proof.\qed
\end{proof}

%-----------------------------------------------------------------
\subsection{Auditability of Hedge-Ratio Computation}
\label{ssec:hedge-audit}
%-----------------------------------------------------------------

\begin{remark}[Operational Requirements for \texttt{HedgeOptimizationOperator}]
\label{rem:hedge-operator}
The \texttt{HedgeOptimizationOperator} must satisfy the following
requirements at every invocation time $t$.
\begin{enumerate}
  \item \textbf{Assumption validation.}
        Before computing $h^*$, the operator must verify
        \textbf{(A1)}--\textbf{(A5)} explicitly:
        \begin{enumerate}
          \item Confirm finite sample second moments
                ($\hat\sigma_V^2 < \infty$, $\hat\sigma_F^2 < \infty$,
                $|\widehat{\Cov}| < \infty$).
          \item Assert $\hat\sigma_F^2 > \epsilon_F$ for a
                pre-configured floor $\epsilon_F > 0$; raise
                \texttt{DegenerateHedgeInstrumentException} if violated.
          \item Assert that no future-dated data timestamps appear in
                the estimation window $[t-W,\,t]$.
        \end{enumerate}

  \item \textbf{Regime-conditioned estimation.}
        Correlation and variance estimates must be obtained from the
        \texttt{MarketRegimeOperator} conditioned on the current
        market regime $s_t \in \mathcal{S}$.
        The regime label $s_t$ and its posterior probability must be
        logged alongside $\hat{h}^*$.

  \item \textbf{Dynamic recomputation.}
        The ratio $\hat{h}^*_t = \widehat{\Cov}_t / \hat\sigma_{F,t}^2$
        must be recomputed at every decision epoch $t$; stale values
        from a prior regime are not permissible.

  \item \textbf{Audit log fields.}
        Each invocation must emit a structured log record containing:
        \begin{enumerate}
          \item Timestamp $t$ and regime label $s_t$.
          \item Estimation window $[t-W,\,t]$ and window length $W$.
          \item Estimates $\hat\sigma_{V,t}^2$, $\hat\sigma_{F,t}^2$,
                $\widehat{\Cov}_t$, $\hat\rho_{VF,t}$.
          \item Computed $\hat{h}^*_t$ and hedge effectiveness
                $\hat\eta_t = \hat\rho_{VF,t}^2$.
          \item Pass/fail status for each of
                \textbf{(A1)}--\textbf{(A5)}.
          \item Any exception raised and the resulting trading action
                (suppression or continuation).
        \end{enumerate}

  \item \textbf{Monotonicity check.}
        If $|\hat{h}^*_t - \hat{h}^*_{t-1}| > \Delta_{\max}$ for a
        pre-configured jump threshold $\Delta_{\max}$, the operator
        must flag the update for human review before applying it.

  \item \textbf{Non-negativity of effectiveness.}
        The operator must assert $\hat\eta_t \geq 0$ after every
        computation; a negative value indicates a numerical error and
        must trigger an alert.
\end{enumerate}
\end{remark}

% ==============================================================
\chapter{Temporal-Leakage Invariant}
\label{ch:leakage}
% ==============================================================

\section{Foundational Definitions}

\begin{definition}[Probability Space]
\label{def:prob-space}
Fix a complete probability space $(\Omega, \mathcal{F}, \mathbb{P})$
throughout this chapter. All random variables and $\sigma$-algebras
are defined with respect to this space.
\end{definition}

\begin{definition}[Filtration]
\label{def:filtration}
A \emph{filtration} is a family
$\mathbb{F} = \{\mathcal{F}_t\}_{t \in \mathcal{T}}$
of sub-$\sigma$-algebras of $\mathcal{F}$ satisfying:
\begin{enumerate}
  \item \textbf{(Monotonicity)}
        $s \leq t \;\Rightarrow\; \mathcal{F}_s \subseteq \mathcal{F}_t$.
  \item \textbf{(Completeness)}
        $\mathcal{F}_0$ contains all $\mathbb{P}$-null sets of
        $\mathcal{F}$.
  \item \textbf{(Right-continuity)}
        $\mathcal{F}_t = \bigcap_{u > t} \mathcal{F}_u$ for all
        $t \in \mathcal{T}$.
\end{enumerate}
The index set $\mathcal{T}$ is either $\mathbb{N}_0$ (discrete time)
or $[0,T]$ (continuous time); both cases are covered by the results
below.
\end{definition}

\begin{definition}[Information Set]
\label{def:info-set}
The \emph{information set} at time $t$, denoted $\calI_t$, is the
sub-$\sigma$-algebra of $\mathcal{F}$ defined by
\begin{equation}
  \calI_t \;:=\; \sigma\!\bigl(X_s : s \leq t\bigr),
  \label{eq:info-set}
\end{equation}
where $\{X_s\}_{s \in \mathcal{T}}$ is the family of all observable
random variables (prices, volumes, macro releases, model outputs,
etc.) whose \emph{realisation} is known no later than calendar time
$t$.  By construction, $\{\calI_t\}_{t \in \mathcal{T}}$ is a
filtration in the sense of Definition~\ref{def:filtration}.
\end{definition}

\begin{remark}[Scope of $\calI_t$]
\label{rem:info-scope}
The following data classes are explicitly \emph{included} in
$\calI_t$:
\begin{enumerate}
  \item All trade prices and bid-ask quotes up to and including
        time $t$.
  \item Macro-economic releases whose official publication timestamp
        is $\leq t$.
  \item Model parameters estimated solely on the sub-sample
        $(-\infty, t]$.
  \item Any derived features (returns, rolling statistics, etc.)
        computed exclusively from the above.
\end{enumerate}
The following data classes are explicitly \emph{excluded} from
$\calI_t$:
\begin{enumerate}
  \item Prices, quotes, or order-book snapshots with timestamps
        $> t$.
  \item Revised macro data published after $t$ (e.g., GDP
        revisions).
  \item Normalisation parameters (means, standard deviations,
        quantile boundaries) estimated on any window that overlaps
        $(t, \infty)$.
  \item Train/test splits that include test data in the training
        set (look-ahead contamination).
\end{enumerate}
\end{remark}

\section{Temporal Validity}

\begin{definition}[Decision Function]
\label{def:decision-fn}
A \emph{decision function} at time $t$ is a map
\begin{align}
  \delta_t \;:\; \Omega &\;\longrightarrow\; \calA,
  \label{eq:decision-fn}
\end{align}
where $\calA$ is the action space (e.g., position sizes, order
types). The decision function $\delta_t$ is
\emph{$\calI_t$-measurable} if
\begin{equation}
  \delta_t^{-1}(B) \;\in\; \calI_t
  \quad \forall\, B \in \mathcal{B}(\calA),
  \label{eq:measurability}
\end{equation}
where $\mathcal{B}(\calA)$ is the Borel $\sigma$-algebra on
$\calA$.
\end{definition}

\begin{definition}[Strategy]
\label{def:strategy}
A \emph{trading strategy} is a sequence (or process)
$\pi = \{\delta_t\}_{t \in \mathcal{T}}$ of decision functions.
\end{definition}

\begin{definition}[Temporal Validity]
\label{def:temporal-validity}
A strategy $\pi = \{\delta_t\}_{t \in \mathcal{T}}$ is
\emph{temporally valid} if and only if every decision function is
measurable with respect to the contemporaneous information set:
\begin{equation}
  \forall\, t \in \mathcal{T}:\quad
  \delta_t \;\text{is } \calI_t\text{-measurable},
  \label{eq:temporal-valid}
\end{equation}
equivalently,
\begin{equation}
  \forall\, t \in \mathcal{T}:\quad
  \mathrm{Decision}_t(\pi) \;\in\; \sigma(\calI_t).
  \label{eq:temporal-valid-alt}
\end{equation}
Conditions~\eqref{eq:temporal-valid} and~\eqref{eq:temporal-valid-alt}
are identical; both forms are stated for cross-reference compatibility.
\end{definition}

\section{Leakage Taxonomy}

\begin{axiom}[Temporal Leakage Violation]
\label{ax:leakage}
A strategy $\pi$ is \emph{temporally invalid} (exhibits
\emph{temporal leakage}) if there exist $t \in \mathcal{T}$ and
$\tau > 0$ with $t + \tau \in \mathcal{T}$ such that
$\delta_t$ is \emph{not} $\calI_t$-measurable, i.e.,
\begin{equation}
  \exists\, B \in \mathcal{B}(\calA) :
  \quad \delta_t^{-1}(B) \;\notin\; \calI_t,
  \quad\text{yet}\quad
  \delta_t^{-1}(B) \;\in\; \calI_{t+\tau}.
  \label{eq:leakage-def}
\end{equation}
Equivalently, $\delta_t$ depends on information in
$\calI_{t+\tau} \setminus \calI_t$.
\end{axiom}

\begin{remark}[Canonical Leakage Sources]
\label{rem:leakage-sources}
The following constitute the exhaustive list of known leakage
mechanisms, each mapped to its formal characterisation
under Axiom~\ref{ax:leakage}:
\begin{enumerate}
  \item \textbf{Future price look-ahead.}
        $\delta_t$ uses $X_{t+\tau}$ (e.g., tomorrow's close)
        directly in the signal computation.
  \item \textbf{Revised macro data.}
        A macro release $M_t$ is replaced by its revised value
        $\tilde{M}_t$ published at $t + \tau$, so $\delta_t$
        effectively conditions on $\tilde{M}_t \in
        \calI_{t+\tau} \setminus \calI_t$.
  \item \textbf{Future news and events.}
        Corporate announcements, earnings surprises, or regulatory
        decisions occurring at $t + \tau$ are used to filter
        training data that precedes $t$.
  \item \textbf{Full-sample normalisation.}
        Normalisation statistics $\mu, \sigma$ are computed over
        the full dataset $[0, T]$; their values depend on
        $\{X_s\}_{s > t}$, placing them in
        $\calI_T \setminus \calI_t$.
  \item \textbf{Train/test contamination.}
        Training samples include observations from the test
        period $[t_{\mathrm{test}}, T]$, causing model parameters
        to be functions of future realisations.
  \item \textbf{Target encoding with future labels.}
        Categorical features are encoded using statistics that
        include labels from periods after $t$.
\end{enumerate}
\end{remark}

\section{Main Results}

\begin{lemma}[Measurability is Hereditary]
\label{lem:measurability-hereditary}
Let $f : \Omega \to \mathbb{R}$ be $\calI_s$-measurable for some
$s \leq t$.  Then $f$ is also $\calI_t$-measurable.
\end{lemma}

\begin{proof}
We verify every step with full rigour.  No assumption is left
implicit; boundary and edge cases are addressed explicitly.

\medskip
\noindent\textbf{Step 1: Validate the domain and type of $f$.}
\begin{enumerate}
  \item By hypothesis, $f : \Omega \to \mathbb{R}$ is a
        well-defined function on the sample space $\Omega$ of the
        probability space $(\Omega, \mathcal{F}, \mathbb{P})$
        introduced in Definition~\ref{def:prob-space}.
  \item The codomain $\mathbb{R}$ is equipped with its Borel
        $\sigma$-algebra $\mathcal{B}(\mathbb{R})$, generated by
        all open intervals $(a, b)$ with $a, b \in \mathbb{R}$,
        $a < b$.
  \item The statement ``$f$ is $\calI_s$-measurable'' means:
        \begin{align}
          \forall\, B \in \mathcal{B}(\mathbb{R}),
          \quad
          f^{-1}(B)
          \;:=\;
          \{\,\omega \in \Omega : f(\omega) \in B\,\}
          \;\in\; \calI_s.
        \end{align}
        This is the standard definition of measurability of a
        real-valued function with respect to a sub-$\sigma$-algebra
        $\calI_s \subseteq \mathcal{F}$.
\end{enumerate}

\medskip
\noindent\textbf{Step 2: Validate the time indices and their
ordering.}
\begin{enumerate}
  \item The index $s$ satisfies $s \in \mathcal{T}$ and the
        hypothesis imposes $s \leq t$ for a given
        $t \in \mathcal{T}$.
  \item \textbf{Boundary case $s = t$.}  If $s = t$ then
        $\calI_s = \calI_t$ by the definition of the filtration,
        and the conclusion $\calI_s \subseteq \calI_t$ holds
        trivially with equality.  The remainder of the proof
        covers the strict case $s < t$ without loss of generality.
  \item \textbf{Boundary case $s = 0$.}  The smallest element of
        $\mathcal{T}$ yields the smallest sub-$\sigma$-algebra
        $\calI_0$.  Definition~\ref{def:filtration}(1) guarantees
        $\calI_0 \subseteq \calI_t$ for all $t \geq 0$, so the
        argument below applies uniformly.
\end{enumerate}

\medskip
\noindent\textbf{Step 3: Apply filtration monotonicity.}
\begin{enumerate}
  \item By Definition~\ref{def:filtration}(1), the family
        $\{\calI_t\}_{t \in \mathcal{T}}$ is a filtration, meaning:
        \begin{align}
          s \leq t
          \;\Longrightarrow\;
          \calI_s \;\subseteq\; \calI_t.
        \end{align}
  \item Since $s \leq t$ holds by hypothesis (Step 2), we
        conclude:
        \begin{align}
          \calI_s \;\subseteq\; \calI_t.
          \label{eq:filtration-mono-hereditary}
        \end{align}
  \item The inclusion~\eqref{eq:filtration-mono-hereditary} is
        between $\sigma$-algebras, i.e., every set $A \in \calI_s$
        also satisfies $A \in \calI_t$.
\end{enumerate}

\medskip
\noindent\textbf{Step 4: Transfer measurability from
$\calI_s$ to $\calI_t$.}
\begin{enumerate}
  \item Let $B \in \mathcal{B}(\mathbb{R})$ be arbitrary.
  \item From Step 1, $f^{-1}(B) \in \calI_s$.
  \item From Step 3, $\calI_s \subseteq \calI_t$, so every
        element of $\calI_s$ belongs to $\calI_t$:
        \begin{align}
          f^{-1}(B) \;\in\; \calI_s
          \;\subseteq\; \calI_t
          \;\Longrightarrow\;
          f^{-1}(B) \;\in\; \calI_t.
        \end{align}
  \item Since $B \in \mathcal{B}(\mathbb{R})$ was arbitrary,
        the implication holds for \emph{all} Borel sets:
        \begin{align}
          \forall\, B \in \mathcal{B}(\mathbb{R}),
          \quad
          f^{-1}(B) \;\in\; \calI_t.
        \end{align}
  \item This is precisely the definition of $f$ being
        $\calI_t$-measurable.
\end{enumerate}

\medskip
\noindent\textbf{Step 5: Conclusion.}
\begin{enumerate}
  \item All hypotheses have been validated in Steps 1--2.
  \item Filtration monotonicity was invoked in Step 3 with its
        precondition $s \leq t$ confirmed.
  \item The measurability transfer in Step 4 is valid for every
        $B \in \mathcal{B}(\mathbb{R})$ without exception,
        covering all boundary and edge cases.
  \item Therefore $f$ is $\calI_t$-measurable, as claimed.\qed
\end{enumerate}
\end{proof}

\begin{lemma}[Future Dependence Implies Non-Measurability]
\label{lem:future-dep}
Let $(\Omega, \mathcal{F}, \mathbb{P})$ be a complete probability
space with filtration $\{\calI_t\}_{t \in \mathcal{T}}$ satisfying
Definition~\ref{def:filtration}.  Let
$g : \Omega \to \mathbb{R}$ be $\calI_{t+\tau}$-measurable for
some $t \in \mathcal{T}$ and $\tau > 0$, and suppose $g$ is
\emph{not} $\calI_t$-measurable.  Then there exists no
$\calI_t$-measurable function $h : \Omega \to \mathbb{R}$
satisfying $h = g$ $\mathbb{P}$-a.s.
\end{lemma}

\begin{proof}
We proceed by contradiction, validating every step and every
edge case explicitly.

\medskip
\noindent\textbf{Step 1: Record all hypotheses and validate their
consistency.}
\begin{enumerate}
  \item \textbf{H1.} $(\Omega, \mathcal{F}, \mathbb{P})$ is a
        complete probability space, meaning: for every
        $N \in \mathcal{F}$ with $\mathbb{P}(N) = 0$ and every
        $A \subseteq N$, we have $A \in \mathcal{F}$.
  \item \textbf{H2.} $\calI_t \subseteq \mathcal{F}$ is a
        sub-$\sigma$-algebra in the filtration, and by
        Definition~\ref{def:filtration}(2), $\calI_t$ contains
        all $\mathbb{P}$-null sets:
        \begin{align}
          \forall\, A \in \mathcal{F},\;
          \mathbb{P}(A) = 0
          \;\Longrightarrow\;
          A \in \calI_t.
        \end{align}
  \item \textbf{H3.} $g$ is $\calI_{t+\tau}$-measurable for
        some $\tau > 0$.
  \item \textbf{H4.} $g$ is \emph{not} $\calI_t$-measurable,
        i.e., there exists at least one Borel set
        $B^* \in \mathcal{B}(\mathbb{R})$ such that
        $g^{-1}(B^*) \notin \calI_t$.
  \item \textbf{Consistency check.}  Since $\tau > 0$, we have
        $t < t + \tau$, so by filtration monotonicity
        $\calI_t \subsetneq \calI_{t+\tau}$ (strict inclusion is
        consistent with H4).  No contradiction arises among H1--H4.
\end{enumerate}

\medskip
\noindent\textbf{Step 2: Introduce the contradictory assumption.}
\begin{enumerate}
  \item Suppose, for the sake of contradiction, that there exists
        a function $h : \Omega \to \mathbb{R}$ such that:
        \begin{enumerate}
          \item[(a)] $h$ is $\calI_t$-measurable, and
          \item[(b)] $h(\omega) = g(\omega)$ for
                $\mathbb{P}$-almost every $\omega \in \Omega$,
                i.e.,
                \begin{align}
                  \mathbb{P}\!\left(\,
                    \{\,\omega \in \Omega :
                    h(\omega) \neq g(\omega)\,\}
                  \,\right)
                  \;=\; 0.
                \end{align}
        \end{enumerate}
  \item Denote the exceptional set by
        $N := \{\omega \in \Omega : h(\omega) \neq g(\omega)\}$,
        so $\mathbb{P}(N) = 0$.
\end{enumerate}

\medskip
\noindent\textbf{Step 3: Invoke completeness to transfer
measurability.}
\begin{enumerate}
  \item We use the following standard result from measure theory:

        \emph{If $(\Omega, \mathcal{F}, \mathbb{P})$ is complete,
        $\calI_t$ contains all $\mathbb{P}$-null sets in
        $\mathcal{F}$, $h$ is $\calI_t$-measurable, and
        $h = g$ $\mathbb{P}$-a.s., then $g$ is
        $\calI_t$-measurable.}

  \item We verify this result step by step.  Let
        $B \in \mathcal{B}(\mathbb{R})$ be arbitrary.  We must
        show $g^{-1}(B) \in \calI_t$.

  \item Decompose the preimage of $B$ under $g$ as follows:
        \begin{align}
          g^{-1}(B)
          &\;=\;
          \bigl[g^{-1}(B) \cap N^c\bigr]
          \;\cup\;
          \bigl[g^{-1}(B) \cap N\bigr],
        \end{align}
        where $N^c = \Omega \setminus N$.

  \item \textbf{First term.}
        On $N^c$, $h(\omega) = g(\omega)$ by definition of $N$.
        Therefore:
        \begin{align}
          g^{-1}(B) \cap N^c
          &\;=\;
          \{\omega \in N^c : g(\omega) \in B\}
          \\[4pt]
          &\;=\;
          \{\omega \in N^c : h(\omega) \in B\}
          \\[4pt]
          &\;=\;
          h^{-1}(B) \cap N^c.
        \end{align}
        Since $h$ is $\calI_t$-measurable, $h^{-1}(B) \in \calI_t$.
        Since $\mathbb{P}(N) = 0$, by H2 we have $N \in \calI_t$,
        hence $N^c \in \calI_t$ (as $\calI_t$ is a
        $\sigma$-algebra).  Therefore:
        \begin{align}
          h^{-1}(B) \cap N^c \;\in\; \calI_t.
        \end{align}

  \item \textbf{Second term.}
        $g^{-1}(B) \cap N \subseteq N$, and
        $\mathbb{P}(N) = 0$.  By H1 (completeness), every subset
        of a null set belongs to $\mathcal{F}$.  By H2,
        $\calI_t$ contains all $\mathbb{P}$-null sets in
        $\mathcal{F}$, so $N \in \calI_t$, and every subset of
        $N$ is a $\mathbb{P}$-null set in $\mathcal{F}$, hence:
        \begin{align}
          g^{-1}(B) \cap N \;\in\; \calI_t.
        \end{align}

  \item \textbf{Union.}
        Since $\calI_t$ is a $\sigma$-algebra and both terms
        belong to $\calI_t$:
        \begin{align}
          g^{-1}(B)
          \;=\;
          \bigl[h^{-1}(B) \cap N^c\bigr]
          \;\cup\;
          \bigl[g^{-1}(B) \cap N\bigr]
          \;\in\; \calI_t.
        \end{align}

  \item Since $B \in \mathcal{B}(\mathbb{R})$ was arbitrary, we
        conclude:
        \begin{align}
          \forall\, B \in \mathcal{B}(\mathbb{R}),
          \quad
          g^{-1}(B) \;\in\; \calI_t,
        \end{align}
        i.e., $g$ is $\calI_t$-measurable.
\end{enumerate}

\medskip
\noindent\textbf{Step 4: Derive the contradiction.}
\begin{enumerate}
  \item Step 3 established that, under the contradictory
        assumption of Step 2, $g$ is $\calI_t$-measurable.
  \item This directly contradicts hypothesis H4, which states
        that $g$ is \emph{not} $\calI_t$-measurable.
  \item The contradiction is strict: there is no escape via
        edge cases, since the argument in Step 3 holds for
        \emph{every} $B \in \mathcal{B}(\mathbb{R})$ and relies
        only on H1--H2, both of which are given by hypothesis.
\end{enumerate}

\medskip
\noindent\textbf{Step 5: Conclusion.}
\begin{enumerate}
  \item The contradictory assumption of Step 2 is false.
  \item Therefore, no $\calI_t$-measurable function
        $h : \Omega \to \mathbb{R}$ satisfies
        $h = g$ $\mathbb{P}$-a.s.
  \item All hypotheses H1--H4 have been used; no step was
        skipped; all boundary conditions ($\tau > 0$,
        completeness, null-set containment) were validated
        explicitly.\qed
\end{enumerate}
\end{proof}

\begin{theorem}[Leakage Implies Invalidity]
\label{thm:leakage-invalidity}
Let $\pi = \{\delta_t\}$ be a strategy that is temporally invalid
in the sense of Axiom~\ref{ax:leakage}.  Then:
\begin{enumerate}
  \item $\pi$ violates the temporal-validity
        condition~\eqref{eq:temporal-valid}.
  \item The profit-and-loss (P\&L) stream generated by $\pi$ in
        backtest \emph{cannot} be replicated in real-time deployment.
  \item Any reported backtest profitability of $\pi$ is
        \emph{vacuous}: it constitutes no evidence of genuine
        predictive skill.
\end{enumerate}
\end{theorem}

\begin{proof}
We establish each claim in sequence.

\medskip
\noindent\textbf{Claim 1: Violation of temporal validity.}
\begin{enumerate}
  \item By hypothesis (Axiom~\ref{ax:leakage}), there exist
        $t \in \mathcal{T}$ and $\tau > 0$ such that
        $\delta_t^{-1}(B) \notin \calI_t$ for some
        $B \in \mathcal{B}(\calA)$.
  \item By Definition~\ref{def:measurability}
        (equation~\eqref{eq:measurability}), $\delta_t$ is
        \emph{not} $\calI_t$-measurable.
  \item By Definition~\ref{def:temporal-validity}
        (condition~\eqref{eq:temporal-valid}), $\pi$ is not
        temporally valid. $\checkmark$
\end{enumerate}

\medskip
\noindent\textbf{Claim 2: Non-replicability of P\&L.}
\begin{enumerate}
  \item Let $\mathrm{PnL}_t(\pi) = \delta_t \cdot (X_{t+1} - X_t)$
        denote the one-step P\&L contribution (generalisation to
        multi-asset or non-linear payoffs is immediate by linearity
        of expectation and does not affect measurability arguments).
  \item Since $\delta_t$ is not $\calI_t$-measurable (established in
        Claim~1), $\delta_t$ cannot be computed from data available at
        time $t$ alone.
  \item In live deployment, only $\calI_t$ is accessible at decision
        time $t$; the set $\calI_{t+\tau} \setminus \calI_t$ is
        strictly unavailable.
  \item By Lemma~\ref{lem:future-dep}, no $\calI_t$-measurable
        substitute $\tilde{\delta}_t$ satisfies
        $\tilde{\delta}_t = \delta_t$ $\mathbb{P}$-a.s.
  \item Therefore $\mathrm{PnL}_t(\pi)$ generated in backtest
        cannot be matched in real time. $\checkmark$
\end{enumerate}

\medskip
\noindent\textbf{Claim 3: Vacuousness of reported profitability.}
\begin{enumerate}
  \item Define the backtest cumulative P\&L:
        \begin{align}
          \Pi^{\mathrm{BT}}(\pi)
          &\;=\; \sum_{t \in \mathcal{T}} \mathrm{PnL}_t(\pi).
          \label{eq:bt-pnl}
        \end{align}
  \item By Claim~2, $\Pi^{\mathrm{BT}}(\pi)$ is computed using
        information $\calI_{t+\tau} \setminus \calI_t$, which is
        unavailable to any real investor at time $t$.
  \item Formally, let $\mathbb{E}^{\calI_t}[\,\cdot\,]$ denote the
        conditional expectation given $\calI_t$.  Predictive skill
        requires:
        \begin{align}
          \mathbb{E}^{\calI_t}\!\bigl[\mathrm{PnL}_t(\pi)\bigr]
          &\;>\; 0
          \quad \forall\, t,
          \label{eq:genuine-skill}
        \end{align}
        computed using only $\calI_t$-measurable $\delta_t$.
  \item Since $\delta_t$ is not $\calI_t$-measurable, the left-hand
        side of~\eqref{eq:genuine-skill} is \emph{undefined} as a
        statement about real-time performance; positive reported
        P\&L reflects exploitation of future information, not skill.
  \item Therefore $\Pi^{\mathrm{BT}}(\pi) > 0$ constitutes no
        evidence of predictive skill. $\checkmark$
\end{enumerate}
All three claims are established; the theorem follows.\qed
\end{proof}

\section{Boundary and Edge Cases}

\begin{proposition}[Zero-Lag Boundary]
\label{prop:zero-lag}
The case $\tau = 0$ is excluded from Axiom~\ref{ax:leakage} by
definition.  A decision function using only data available at exactly
time $t$ (i.e., $\delta_t$ is $\calI_t$-measurable) does \emph{not}
constitute leakage, even if the realisation of $X_t$ arrives
at time $t$ with zero latency.
\end{proposition}

\begin{proof}
If $\tau = 0$, then $\calI_{t+\tau} = \calI_t$, so
$\calI_{t+\tau} \setminus \calI_t = \emptyset$.  Condition
\eqref{eq:leakage-def} cannot be satisfied, and the strategy is
temporally valid by Definition~\ref{def:temporal-validity}.\qed
\end{proof}

\begin{proposition}[Infinitesimal Look-Ahead]
\label{prop:infinitesimal}
In continuous time, for any $\tau > 0$ (however small),
temporal leakage is present if $\delta_t$ depends on
$X_{t+\tau}$.  There is no ``negligible'' leakage threshold;
any strictly positive $\tau$ violates
condition~\eqref{eq:temporal-valid}.
\end{proposition}

\begin{proof}
\begin{enumerate}
  \item Fix any $\tau > 0$.  By right-continuity of the filtration
        (Definition~\ref{def:filtration}(3)),
        $\calI_t = \bigcap_{u > t} \calI_u$.
  \item For any $\epsilon \in (0, \tau)$, $X_{t+\epsilon}$ is
        $\calI_{t+\epsilon}$-measurable but need not be
        $\calI_t$-measurable if the process $\{X_s\}$ is not
        $\mathbb{F}$-predictable.
  \item In particular, if $X$ has independent increments (e.g., is
        a martingale), $X_{t+\tau} - X_t$ is independent of
        $\calI_t$ for all $\tau > 0$, so knowledge of $X_{t+\tau}$
        strictly enlarges the information available beyond
        $\calI_t$.
  \item Therefore any dependence on $X_{t+\tau}$ for $\tau > 0$
        constitutes a violation of~\eqref{eq:temporal-valid}.\qed
\end{enumerate}
\end{proof}

\begin{proposition}[Leakage Transitivity]
\label{prop:leakage-transitivity}
If feature $F_t$ is a deterministic function of a leaked variable
$Z_{t+\tau}$ (with $\tau > 0$), then any strategy using $F_t$
inherits the leakage of $Z_{t+\tau}$.
\end{proposition}

\begin{proof}
\begin{enumerate}
  \item Let $F_t = \phi(Z_{t+\tau})$ for some measurable function
        $\phi$.
  \item Since $Z_{t+\tau} \in \calI_{t+\tau} \setminus \calI_t$
        (by hypothesis), $\sigma(Z_{t+\tau}) \subseteq \calI_{t+\tau}$
        but $\sigma(Z_{t+\tau}) \not\subseteq \calI_t$.
  \item As $F_t$ is $\sigma(Z_{t+\tau})$-measurable (being a
        measurable function of $Z_{t+\tau}$), $F_t$ is
        $\calI_{t+\tau}$-measurable but not $\calI_t$-measurable.
  \item Therefore any $\delta_t$ that is a measurable function of
        $F_t$ is not $\calI_t$-measurable, and the strategy violates
        condition~\eqref{eq:temporal-valid}.\qed
\end{enumerate}
\end{proof}

\section{Operational Enforcement}

\begin{remark}[CI Enforcement Gate]
\label{rem:ci-gate}
A \texttt{TemporalLeakageDetector} continuous-integration (CI) gate
must enforce condition~\eqref{eq:temporal-valid} on every merge to
the \texttt{main} branch.  The gate must satisfy the following
requirements:
\begin{enumerate}
  \item \textbf{Timestamp audit.}
        Every data source must carry a verified ingestion timestamp
        $\hat{t} \leq t$; any source with $\hat{t} > t$ must cause
        the build to fail.
  \item \textbf{Normalisation-window check.}
        All rolling statistics (means, variances, quantiles) must be
        computed over expanding or strictly past-anchored windows;
        any window boundary exceeding the current bar must be flagged.
  \item \textbf{Feature-lineage graph.}
        A directed acyclic graph (DAG) of feature dependencies must
        be constructed at build time; any path from a node
        $X_{t+\tau}$ ($\tau > 0$) to a decision node $\delta_t$
        must cause the build to fail (Proposition~%
        \ref{prop:leakage-transitivity}).
  \item \textbf{Train/test boundary enforcement.}
        The training index set $\mathcal{T}_{\mathrm{train}}$ must
        satisfy $\max(\mathcal{T}_{\mathrm{train}}) <
        \min(\mathcal{T}_{\mathrm{test}})$; any overlap must cause
        the build to fail.
  \item \textbf{Audit log.}
        The gate must emit a machine-readable audit record containing:
        \begin{enumerate}
          \item The commit hash and timestamp.
          \item The complete feature-lineage DAG.
          \item Pass/fail status for each of the four checks above.
          \item Any violation details sufficient for third-party
                reproduction.
        \end{enumerate}
\end{enumerate}
\end{remark}

% ==============================================================
\chapter{Operator Composition Algebra}
\label{ch:operators}
% ==============================================================

\section{Operator Definition}

\begin{definition}[Operator]
\label{def:operator}
Let $\calO$ be the set of all operators.  Each operator
$O_i \in \calO$ is a pure function
\begin{equation}
  O_i : \calS \times \calI_t \;\to\; \calS',
  \label{eq:operator-def}
\end{equation}
where $\calS$ is the shared state space and $\calI_t$ the information
set at time $t$.
\end{definition}

\section{Sequential Composition}

\begin{definition}[Sequential Composition]
\label{def:seq-composition}
Given $O_1, O_2 \in \calO$, sequential composition is
\begin{equation}
  (O_1 \circ O_2)(s, \calI_t) \;=\; O_1\!\left(O_2(s, \calI_t),\, \calI_t\right).
  \label{eq:seq-comp}
\end{equation}
\end{definition}

\begin{proposition}[Associativity of Composition]
\label{prop:assoc}
Operator composition is associative:
$(O_1 \circ O_2) \circ O_3 = O_1 \circ (O_2 \circ O_3)$.
\end{proposition}

\begin{proof}
For any $(s, \calI_t)$:
\begin{align*}
  ((O_1 \circ O_2) \circ O_3)(s, \calI_t)
  &= (O_1 \circ O_2)(O_3(s, \calI_t), \calI_t)\\
  &= O_1(O_2(O_3(s, \calI_t), \calI_t), \calI_t)\\
  &= (O_1 \circ (O_2 \circ O_3))(s, \calI_t).
\end{align*}\qed
\end{proof}

\begin{axiom}[Pipeline Temporal Invariant]
\label{ax:pipeline}
A valid operator pipeline $\Pi = O_n \circ \cdots \circ O_1$ must satisfy
\begin{equation}
  \Pi(s, \calI_t) \;\in\; \sigma(\calI_t)
  \quad \forall\, s \in \calS,\; t \in \calT.
  \label{eq:pipeline-invariant}
\end{equation}
\end{axiom}

\section{Agent as Weighted Ensemble}

\begin{definition}[Agent]
\label{def:agent}
An \emph{agent} $A$ is a tuple $(O_1, \ldots, O_k, w)$ where $O_i$ are
operators and $w \in \Delta^{k-1}$ is a weight vector on the probability
simplex
\begin{equation}
  \Delta^{k-1} = \left\{w \in \R^k \;\middle|\;
    \sum_{i=1}^{k} w_i = 1,\; w_i \geq 0\right\}.
\end{equation}
The \emph{agent signal} is
\begin{equation}
  \hat{y}_t \;=\; \sum_{i=1}^{k} w_i\,O_i(s, \calI_t).
  \label{eq:agent-signal}
\end{equation}
\end{definition}

\begin{proposition}[Ensemble Variance Bound]
\label{prop:ensemble-var}
Under pairwise independence among operators with forecast variances
$\sigma_i^2$:
\begin{equation}
  \Var(\hat{y}_t) \;=\; \sum_{i=1}^{k} w_i^2\,\sigma_i^2
                  \;\leq\; \max_i\,\sigma_i^2.
  \label{eq:ensemble-var}
\end{equation}
\end{proposition}

\begin{proof}
Under independence and unit weights satisfying $\sum w_i = 1$,
$\Var(\hat{y}_t) = \sum w_i^2 \sigma_i^2 \leq (\max_i \sigma_i^2)\sum w_i^2
\leq \max_i \sigma_i^2$ since $\sum w_i^2 \leq (\sum w_i)^2 = 1$.\qed
\end{proof}

% ==============================================================
\chapter{Four-Domain Asset Partition}
\label{ch:four-domain}
% ==============================================================

\begin{definition}[Four-Domain Partition]
\label{def:four-domain}
The full asset universe $\calU$ is partitioned as
\begin{equation}
  \calU \;=\; \mathcal{D}_{E} \cup \mathcal{D}_{A}
            \cup \mathcal{D}_{F} \cup \mathcal{D}_{W},
  \qquad
  \mathcal{D}_i \cap \mathcal{D}_j = \emptyset \;\;(i \neq j),
  \label{eq:four-domain}
\end{equation}
where $\mathcal{D}_E$ (Earth) = Portfolio/Equities,
$\mathcal{D}_A$ (Air) = Rates/Macro Hedging,
$\mathcal{D}_F$ (Fire) = Volatility/Options,
$\mathcal{D}_W$ (Water) = Crypto/Cross-Market Liquidity.
\end{definition}

\begin{definition}[Cross-Domain Influence Matrix]
\label{def:influence-matrix}
The directed influence matrix $\Gamma \in \R^{4\times 4}$ is defined by
$\Gamma_{ij} = \partial\hat{y}_t^{(i)}/\partial\hat{y}_t^{(j)}$ for $i \neq j$.
\end{definition}

\begin{definition}[Regime-Conditioned Allocation]
\label{def:regime-alloc}
For regime $r \in \calR$, the optimal cross-domain allocation vector
$\alpha^{*} \in \Delta^{3}$ solves
\begin{equation}
  \alpha^{*} \;=\; \arg\max_{\alpha \in \Delta^{3}}\;
  \E\!\left[\sum_{i=1}^{4}\alpha_i\,\mu_i^{(r)}
            - \lambda\,\alpha^\top \Sigma^{(r)}\alpha\right],
  \label{eq:regime-alloc}
\end{equation}
where $\mu^{(r)} \in \R^4$ and $\Sigma^{(r)} \in \R^{4\times 4}$ are
regime-conditioned return vector and covariance matrix.
\end{definition}

% ==============================================================
\chapter{Dependency Graph Acyclicity}
\label{ch:dag}
% ==============================================================

\begin{definition}[Operator Dependency Graph]
\label{def:dep-graph}
The \emph{operator dependency graph} $G = (V, E)$ has vertex set
$V = \calO$ and a directed edge $(O_i, O_j) \in E$ if $O_j$ consumes
output produced by $O_i$.
\end{definition}

\begin{axiom}[DAG Invariant]
\label{ax:dag}
$G$ must be a directed acyclic graph (DAG):
\begin{equation}
  \nexists\;\text{cycle}\;
  O_{i_1} \to O_{i_2} \to \cdots \to O_{i_1} \;\text{in } G.
  \label{eq:dag}
\end{equation}
\end{axiom}

\begin{definition}[Topological Execution Order]
\label{def:topo-order}
A valid execution ordering $\tau : V \to \{1, \ldots, |V|\}$ satisfies
\begin{equation}
  (O_i, O_j) \in E \;\Longrightarrow\; \tau(O_i) < \tau(O_j).
  \label{eq:topo-order}
\end{equation}
\end{definition}

\begin{theorem}[Existence of Topological Order]
\label{thm:topo-exists}
If $G$ satisfies the DAG invariant~\eqref{eq:dag}, a topological ordering
$\tau$ satisfying~\eqref{eq:topo-order} exists and can be computed in
$O(|V| + |E|)$ time by Kahn's algorithm.
\end{theorem}

\begin{proof}
Kahn's algorithm iteratively removes source vertices (in-degree zero) and
appends them to the ordering.  If $G$ is a DAG, every non-empty subgraph
has at least one source vertex; hence the algorithm terminates with a
complete ordering.  Each vertex and edge is processed at most once,
giving $O(|V|+|E|)$ time.  If the algorithm terminates before all vertices
are removed, a cycle is detected, violating the DAG invariant.\qed
\end{proof}

% ==============================================================
\chapter{Idempotency and Referential Transparency}
\label{ch:idempotency}
% ==============================================================

\begin{definition}[Referential Transparency]
\label{def:ref-transparency}
An operator $O$ is \emph{referentially transparent} if
\begin{equation}
  s \equiv s' \pmod{\calI_t}
  \;\Longrightarrow\;
  O(s, \calI_t) = O(s', \calI_t).
  \label{eq:ref-trans}
\end{equation}
\end{definition}

\begin{definition}[Idempotency]
\label{def:idempotency}
An operator $O$ is \emph{idempotent} if
\begin{equation}
  O\!\left(O(s, \calI_t),\, \calI_t\right) = O(s, \calI_t)
  \quad \forall\, s \in \calS,\; t \in \calT.
  \label{eq:idempotent}
\end{equation}
\end{definition}

\begin{proposition}[Safe Parallelisation]
\label{prop:safe-parallel}
Idempotent, referentially transparent operators may be safely memoised,
parallelised, and replayed without side effects.
\end{proposition}

\begin{proof}
Referential transparency guarantees that the output depends only on
$\calI_t$, not on mutable external state.  Idempotency guarantees that
multiple applications do not accumulate errors.  Together these properties
make concurrent execution and replay semantically equivalent to single
sequential execution.\qed
\end{proof}

% ==============================================================
\chapter{Risk-Adjusted Objective Function}
\label{ch:objective}
% ==============================================================

\begin{definition}[Risk-Adjusted Objective]
\label{def:objective}
For discount factor $\gamma \in (0,1]$, CVaR penalty $\lambda_1 \geq 0$,
and drawdown penalty $\lambda_2 \geq 0$, the platform maximises
\begin{equation}
  \mathcal{J}(\pi) \;=\;
  \E\!\left[\sum_{t=0}^{T}\gamma^t\,r_t\right]
  - \lambda_1\,\CVaR_\alpha(\pi)
  - \lambda_2\,\MaxDD(\pi).
  \label{eq:objective}
\end{equation}
\end{definition}

\begin{definition}[Feasible Strategy Set]
\label{def:feasible}
The constrained strategy space is
\begin{equation}
  \Pi_F \;=\;
  \left\{\pi \;\middle|\;
    \Var(\pi) \leq \sigma_{\max}^2,\;
    \|\alpha(\pi)\|_1 \leq L_{\max},\;
    \alpha(\pi) \in \Delta^{n-1}
  \right\}.
  \label{eq:feasible}
\end{equation}
\end{definition}

\begin{remark}
Success is not defined as maximising predicted profit.  It is defined as
finding decisions that remain defensible when models disagree, markets
change regime, forecasts fail, correlations break, and uncertainty becomes
materially greater.
\end{remark}

% ==============================================================
\chapter{Type-Safe Language-Layer Assignment}
\label{ch:languages}
% ==============================================================

\begin{definition}[Language Layer Assignment]
\label{def:lang-assignment}
Let $\calL = \{\texttt{Python}, \texttt{C++}, \texttt{Rust}, \texttt{Go}\}$.
The layer assignment function $\ell : \calO \to \calL$ must satisfy
\begin{align}
  O \in \calO_{\text{research}}  &\;\Longrightarrow\; \ell(O) = \texttt{Python}, \label{eq:lang-py}\\
  O \in \calO_{\text{numerics}}  &\;\Longrightarrow\; \ell(O) \in \{\texttt{C++},\,\texttt{Rust}\}, \label{eq:lang-num}\\
  O \in \calO_{\text{infra}}     &\;\Longrightarrow\; \ell(O) \in \{\texttt{Go},\,\texttt{Rust}\}. \label{eq:lang-infra}
\end{align}
\end{definition}

\begin{axiom}[Latency Bound]
\label{ax:latency}
For any operator $O$ in a latency-sensitive layer, the wall-clock
execution time $\tau_O$ must satisfy
\begin{equation}
  \tau_O \;\leq\; \tau_{\max}, \qquad \tau_{\max} = 10^{-3}\,\text{s}.
  \label{eq:latency}
\end{equation}
\end{axiom}

% ==============================================================
\chapter{Regime Detection as a Hidden Markov Model}
\label{ch:hmm}
% ==============================================================

\begin{definition}[HMM Regime Model]
\label{def:hmm}
The regime sequence $\{R_t\}_{t \geq 0}$ follows a first-order Markov
chain with transition matrix $\Pi \in \R^{K\times K}$,
$\sum_j \Pi_{ij} = 1$.  Observed features are drawn from regime-conditioned
Gaussians:
\begin{equation}
  \bfx_t \mid R_t = r_k \;\sim\; \mathcal{N}(\mu_k,\, \Sigma_k).
  \label{eq:hmm-emission}
\end{equation}
\end{definition}

\begin{definition}[Filtered Regime Belief]
\label{def:hmm-belief}
The filtered posterior $\beta_t(k) = P(R_t = r_k \mid \bfx_{1:t})$ is
updated by
\begin{equation}
  \beta_t(k)
  \;=\;
  \frac{p(\bfx_t \mid R_t = r_k)
        \displaystyle\sum_{i=1}^{K}\Pi_{ik}\,\beta_{t-1}(i)}
       {\displaystyle\sum_{j=1}^{K} p(\bfx_t \mid R_t = r_j)
        \displaystyle\sum_{i=1}^{K}\Pi_{ij}\,\beta_{t-1}(i)}.
  \label{eq:hmm-update}
\end{equation}
\end{definition}

\begin{definition}[MAP Regime Estimate]
\label{def:map-regime}
The maximum a posteriori regime estimate at time $t$ is
\begin{equation}
  \hat{R}_t \;=\; \arg\max_{k \in \{1,\ldots,K\}}\;\beta_t(k).
  \label{eq:map-regime}
\end{equation}
\end{definition}

\begin{theorem}[Normalisation of Filtered Beliefs]
\label{thm:hmm-norm}
$\sum_{k=1}^{K}\beta_t(k) = 1$ for all $t$.
\end{theorem}

\begin{proof}
The denominator of~\eqref{eq:hmm-update} is exactly
$\sum_{k=1}^{K}$ of the numerator, so the ratio sums to unity by
construction.  By induction, if $\sum_k \beta_{t-1}(k) = 1$, then the
predictive distribution $\sum_k \sum_i \Pi_{ik}\beta_{t-1}(i) = \sum_i
\beta_{t-1}(i)\sum_k \Pi_{ik} = 1$, and normalisation at each step
preserves the sum-to-one property.\qed
\end{proof}

% ==============================================================
\chapter{Uncertainty Quantification:\\Epistemic--Aleatoric Decomposition}
\label{ch:uncertainty}
% ==============================================================

\begin{definition}[Total Predictive Variance Decomposition]
\label{def:var-decomp}
For forecast target $y_{t+h}$ and model parameters $\theta$:
\begin{equation}
  \Var(y_{t+h} \mid \calI_t)
  \;=\;
  \underbrace{\E_\theta\!\left[\Var(y_{t+h} \mid \theta, \calI_t)\right]}_{\text{aleatoric}}
  \;+\;
  \underbrace{\Var_\theta\!\left[\E(y_{t+h} \mid \theta, \calI_t)\right]}_{\text{epistemic}}.
  \label{eq:var-decomp}
\end{equation}
\end{definition}

\begin{definition}[Normalised Epistemic Fraction]
\label{def:epistemic-fraction}
\begin{equation}
  \rho_t \;=\;
  \frac{\Var_\theta\!\left[\E(y_{t+h} \mid \theta, \calI_t)\right]}
       {\Var(y_{t+h} \mid \calI_t)}.
  \label{eq:epistemic-frac}
\end{equation}
\end{definition}

\begin{axiom}[Epistemic Signal Gate]
\label{ax:epistemic-gate}
For a configurable epistemic threshold $\rho_{\max} \in (0, 1)$, the
gated signal is
\begin{equation}
  \tilde{s}_t \;=\; s_t \cdot \mathbf{1}\!\left[\rho_t \leq \rho_{\max}\right].
  \label{eq:epistemic-gate}
\end{equation}
Signals generated under excessive model disagreement are suppressed.
\end{axiom}

\begin{proposition}[Boundedness of Epistemic Fraction]
\label{prop:epistemic-bounded}
$\rho_t \in [0, 1]$ for all $t$.
\end{proposition}

\begin{proof}
By the law of total variance~\eqref{eq:var-decomp}, the epistemic term
is a non-negative summand of the total variance, so $0 \leq
\Var_\theta[\E(\cdot)] \leq \Var(y_{t+h} \mid \calI_t)$.  Dividing by the
total variance (assumed positive) gives $\rho_t \in [0, 1]$.\qed
\end{proof}

% ==============================================================
\chapter{Synthesis: Ideal--Real Regime Separation and\\
Na\"{i}ve--Oracle Duality}
\label{ch:synthesis}
% ==============================================================

\section{Ideal and Real Execution Environments}

\begin{definition}[Ideal and Real Execution Environments]
\label{def:ideal-real}
The \emph{ideal execution environment} $\Ideal$ operates on complete,
noise-free information: $\calI_t^{\Ideal} = \sigma\!\left(\bigcup_{\tau \leq t}
\calI_\tau\right)$ with zero uncertainty ($\hat{U}_t = 0$) and zero
slippage.  The \emph{real execution environment} $\Real$ operates on
$\calI_t$ only, with non-zero uncertainty $\hat{U}_t > 0$, transaction
costs, and finite liquidity.
\end{definition}

\begin{definition}[Regime-Separation Gap]
\label{def:gap}
The \emph{regime-separation gap} between $\Ideal$ and $\Real$ under
strategy $\pi$ is
\begin{equation}
  \Gap(\pi) \;=\; \mathcal{J}^{\Ideal}(\pi) - \mathcal{J}^{\Real}(\pi),
  \label{eq:gap}
\end{equation}
where $\mathcal{J}$ is the risk-adjusted objective~\eqref{eq:objective}.
\end{definition}

\begin{theorem}[Non-Negativity of the Gap]
\label{thm:gap-nonneg}
$\Gap(\pi) \geq 0$ for all strategies $\pi \in \Pi_F$.
\end{theorem}

\begin{proof}
The ideal environment is an information-theoretic superset of the real
environment: $\calI_t^{\Ideal} \supseteq \calI_t$.  Any decision rule
feasible in $\Real$ is also feasible in $\Ideal$, so the supremum of
$\mathcal{J}$ over feasible strategies is at least as large in $\Ideal$
as in $\Real$.  Since $\Gap$ compares the same strategy $\pi$ in both
environments and the ideal environment removes all uncertainty-induced
losses, $\mathcal{J}^{\Ideal}(\pi) \geq \mathcal{J}^{\Real}(\pi)$,
giving $\Gap(\pi) \geq 0$.\qed
\end{proof}

\section{Na\"{i}ve--Oracle Duality}

\begin{definition}[Na\"{i}ve and Oracle Strategies]
\label{def:naive-oracle}
The \emph{na\"{i}ve strategy} $\Naive$ applies equal weights across all
operators and never adapts to regime uncertainty.  The \emph{oracle strategy}
$\Oracle$ has access to the true regime $R_t$ at each step and optimises
allocation exactly.
\end{definition}

\begin{definition}[Frontier Model]
\label{def:frontier}
A \emph{frontier model} $\Frontier$ is any strategy satisfying
\begin{equation}
  \mathcal{J}(\Naive) \;\leq\; \mathcal{J}(\Frontier) \;\leq\; \mathcal{J}(\Oracle)
\end{equation}
for all market regimes in $\calR$, with equality on the left when
uncertainty is maximal ($\hat{U}_t = 1$) and equality on the right in
the limit of vanishing uncertainty ($\hat{U}_t \to 0$).
\end{definition}

\begin{conjecture}[Frontier Convergence]
\label{conj:frontier}
As the number of training regimes $|\calR|$ grows and the model capacity
of the \texttt{MarketRegimeOperator} increases, any frontier model
$\Frontier$ satisfying the HMM belief update~\eqref{eq:hmm-update} and
the epistemic gate axiom~\eqref{eq:epistemic-gate} converges in
distribution to the oracle strategy:
\begin{equation}
  \mathcal{J}(\Frontier) \;\to\; \mathcal{J}(\Oracle) \quad
  \text{as } |\calR| \to \infty.
\end{equation}
\end{conjecture}

\begin{remark}
Conjecture~\ref{conj:frontier} is an open research problem.  The formal
framework developed in Chapters~\ref{ch:regime}--\ref{ch:hmm} provides
the rigorous substrate required to make the conjecture precise and to
design empirical tests for it.
\end{remark}

% ==============================================================
\chapter*{References}
\addcontentsline{toc}{chapter}{References}
% ==============================================================

\begin{enumerate}[leftmargin=*, label={[R\arabic*]}]
  \item Markowitz, H.\ (1952). Portfolio Selection. \textit{Journal of Finance}, 7(1), 77--91.
  \item Kelly, J.\ L.\ (1956). A New Interpretation of Information Rate. \textit{Bell System Technical Journal}, 35(4), 917--926.
  \item Sharpe, W.\ F.\ (1966). Mutual Fund Performance. \textit{Journal of Business}, 39(1), 119--138.
  \item Sortino, F.\ A.\ \& van der Meer, R.\ (1991). Downside Risk. \textit{Journal of Portfolio Management}, 17(4), 27--31.
  \item Rockafellar, R.\ T.\ \& Uryasev, S.\ (2000). Optimization of Conditional Value-at-Risk. \textit{Journal of Risk}, 2(3), 21--41.
  \item Hamilton, J.\ D.\ (1989). A New Approach to the Economic Analysis of Nonstationary Time Series and the Business Cycle. \textit{Econometrica}, 57(2), 357--384.
  \item Ederington, L.\ H.\ (1979). The Hedging Performance of the New Futures Markets. \textit{Journal of Finance}, 34(1), 157--170.
  \item Ben-Tal, A.\ \& Nemirovski, A.\ (1998). Robust Convex Optimization. \textit{Mathematics of Operations Research}, 23(4), 769--805.
  \item Shannon, C.\ E.\ (1948). A Mathematical Theory of Communication. \textit{Bell System Technical Journal}, 27(3), 379--423.
  \item Young, T.\ W.\ (1991). Calmar Ratio: A Smoother Tool. \textit{Futures Magazine}, 20(1).
  \item Lopez de Prado, M.\ (2018). \textit{Advances in Financial Machine Learning}. Wiley.
  \item Pardo, R.\ (2008). \textit{The Evaluation and Optimization of Trading Strategies} (2nd ed.). Wiley.
\end{enumerate}

\end{document}